\documentclass{article}
\usepackage[utf8]{inputenc}
\usepackage[T1]{fontenc} % Output font encoding for international characters
\usepackage{mathptmx} %times new roman
\usepackage{mathrsfs} 
\usepackage{dsfont} 
%\usepackage{contour}
 
 %times new roman
%\usepackage[backend=bibtex,style=authoryear,natbib=true, maxcitenames=1]{biblatex} % Use the bibtex backend with the authoryear citation style (which resembles APA)
\usepackage[autostyle=true]{csquotes} % Required to generate language-dependent quotes in the bibliography
%\renewcommand{\chapte a Gaussian stochastic processrmarkformat}{\thechapter\autodot{.}\enspace} %for header
%\usepackage{titlesec}

\usepackage{amsmath,amsfonts,amssymb,amsthm}
\usepackage{graphicx}%
\usepackage{color}
\usepackage[dvipsnames]{xcolor}
\usepackage{ulem}
\usepackage{enumerate}
\usepackage{url}
\usepackage{bm}
\usepackage{bbm}
\usepackage{float}
\usepackage{mathtools}
\usepackage[toc,page]{appendix}
%\usepackage{setspace}
%\usepackage{subcaption}
%\usepackage{lipsum}
\usepackage{needspace}
\usepackage{xpatch}
\usepackage{gensymb}
\usepackage{cancel}
\usepackage{ifthen}
%-------------------------------------------
\usepackage[textsize=scriptsize]{todonotes}
%-------------------------------------------


\xpretocmd{\section}{\needspace{6\baselineskip}}{}{}
\usepackage[section]{placeins}

\newcommand{\B}[1]{{\textcolor{blue}{[Bruno: \textit{#1}]}}}
\newcommand{\BL}[1]{{\textcolor{blue}{[\textcolor{red}{Question => Luc}: \textit{#1})}}}
\newcommand{\Ball}[1]{{\textcolor{blue}{[ \textbf{bruno $\Rightarrow$ \textcolor{red}{All: }}\textit{#1} )}}}

\renewcommand{\L}[1]{{\textcolor{Brown}{[Ludovic: \textit{#1}]}}}
\newcommand{\ludo}{\color{Brown}}
\newcommand{\Luccomment}[1]{\todo[color=yellow]{#1}}
\newcommand{\Lucinline}[1]{\todo[inline,color=yellow]{#1}}
\newcommand{\New}{\color{ForestGreen} }
\newcommand{\Luc}[1]{{\textcolor{magenta}{[Luc: \textit{#1}]}}}
\newcommand{\Chk}{\color{red} }
%\newcommand{\Luc}{\color{magenta} }

\newcommand{\TBC}{\bigskip \B{To be continued} \pagebreak}
\DeclareMathOperator*{\argmin}{\arg\!\min} 
\DeclareMathOperator*{\argmax}{\arg\!\max} 
\DeclareMathAlphabet\mathbfcal{OMS}{cmsy}{b}{n}

%\graphicspath{{Figures/}}

%%% Useful Macros %%%
\newcommand{\Sec}[2]{\section{#1\label{Sec:#2}}}
\newcommand{\sSec}[2]{\subsection{#1 \label{Sec:#2}}}
\newcommand{\ssSec}[2]{\subsubsection{#1 \label{Sec:#2}}}

\newcommand{\setZ}{\m, whileathbbm{Z}}
\newcommand{\setR}{\mathbbm{R}}
\newcommand{\deq}{\stackrel {\rm def}{=}}
\newcommand{\Prior}[1]{{\textcolor{red}{PRIORITE: \textbf{#1}]}}}

\newcommand{\mycaption}[2]{\caption{\textbf{#1}{ #2}}}
%\newcommand{\ang}[2]{\widehat{\pare{\vec{#1},\vec{#2}}}}
\newcommand{\conv}[1]{\stackrel{#1}{\ast}}
\newcommand{\pare}[1]{\left(\, #1 \, \right)}
\newcommand{\brk}[1]{\langle\, #1  \, \rangle}
\newcommand{\scal}[2]{\langle\, #1 | #2 \, \rangle}
\newcommand{\bra}[1]{\left[\, #1 \, \right]}
\newcommand{\abs}[1]{\left|\, #1 \, \right|}
\newcommand{\brac}[2]{\left[\, #1 \, | \, #2 \right]}
\newcommand{\Set}[1]{\left\{\, #1 \, \right\}}
\newcommand{\Setc}[2]{\left\{\, #1 \, ; \, #2 \right\}}
\newcommand{\vect}[1]{\pare{\begin{array}{ccc}#1\end{array}}}
\newcommand{\ve}{\vec{e}}
%\newcommand{\vv}{\vec{v}}
\newcommand{\vv}{\bold{v}}
\newcommand{\vs}{\bold{s}}
\newcommand{\vy}{\bold{y}}
\newcommand{\tvv}{\bold{v}^T}
\newcommand{\vx}{\vec{x}}
\newcommand{\vB}{\bold{B}}
%\newcommand{\vB}{\vec{B}}
\newcommand{\vD}{\bold{D}}
%\newcommand{\vD}{\vec{D}}
\newcommand{\vF}{\bold{F}}
%\newcommand{\vF}{\vec{F}}
\newcommand{\vf}{\bold{f}}
%\newcommand{\vf}{\vec{f}}
\newcommand{\vN}{\bold{N}}
%\newcommand{\vN}{\vec{N}}
\newcommand{\vcF}{\mathbfcal{F}}
%\newcommand{\vcF}{\vec{\cF}}
%\newcommand{\vu}{\vec{u}}
\newcommand{\vu}{\bold{u}}
%\newcommand{\vZ}{\vec{Z}}
\newcommand{\vZ}{\bold{Z}}
%\newcommand{\vX}{\vec{X}}
\newcommand{\vX}{\bold{X}}
\newcommand{\vXT}{\bold{X}^T}
%\newcommand{\vcC}{\vec{\cC}}
\newcommand{\tvw}{\bold{w}^T}
\newcommand{\vw}{\bold{w}}
\newcommand{\vW}{\bold{W}}
\newcommand{\tP}{\tilde{P}}
\newcommand{\tX}{\tilde{X}}
\newcommand{\tY}{\tilde{Y}}
\newcommand{\vcC}{\vec{\cC}}
\newcommand{\vcR}{\vec{\cR}}
\newcommand{\vcX}{\vec{\cX}}
\newcommand{\vcXs}{\vec{\cX}^\ast}
\newcommand{\vgamma}{\vec{\gamma}}
\newcommand{\vnabla}{\vec{\nabla}}
\newcommand{\diag}{\mathop{\mathrm{diag}}}
\newcommand{\lamb}[2]{\lambda^{(#1)}_{#2}}
\newcommand{\Pp}[2]{\cP^{(#1)}_{#2}}
\newcommand{\Ppm}[2]{\pare{\cP^{(#1)}_{#2}}^{-1}}
\newcommand{\CRep}[2]{\varpi_{#1#2}}
\newcommand{\Om}[1]{\Omega^{(#1)}}
\newcommand{\Lc}[1]{\cL^{(#1)}}
\newcommand{\vCc}[1]{\vcC^{(#1)}}
\newcommand{\vXx}[1]{\vcX^{(#1)}}
\newcommand{\vFf}[1]{\vcF^{(#1)}}
\newcommand{\tOm}[2]{t^{(#1)}_{#2}}
\newcommand{\ee}[2]{e_{#1}^{(#2)}}
\newcommand{\Flow}[1]{\Phi^{(#1)}}

\newcommand{\Prob}[1]{\mathds{P}\bra{#1}}
\newcommand{\Exp}[2]{\mathds{E}_{#1}\bra{#2}}
\newcommand{\Cell}[2]{{#1}_{#2}} %Cell type, cell index
\newcommand{\A}[2]{a_{\Cell{#1}{#2}}} %Cell activity
\newcommand{\V}[2]{V_{\Cell{#1}{#2}}} %Voltage Cell type, cell index
\renewcommand{\l}[2]{\nu_{\Cell{#1}{#2}}} %Frequency type, cell index
\newcommand{\Ce}[2]{C_{\Cell{#1}{#2}}} %Capacity type, cell index
\newcommand{\n}[2]{n^{#1}_{#2}} %Capacity type, cell index
\newcommand{\K}[2]{\cK_{\Cell{#1}{#2}}} %Kernel Cell type, cell index
\newcommand{\KS}[2]{\cK_{\Cell{#1}{#2},S}}  %Kernel Cell type, cell index
\newcommand{\KT}[2]{\cK_{\Cell{#1}{#2},T}}  %Kernel Cell type, cell index
\newcommand{\Vd}[2]{V^2_{\Cell{#1}{#2}}}
\newcommand{\E}[3]{\cE^{(#1)}_{\Cell{#2}{#3}}}
\newcommand{\hE}[3]{\hcE^{(#1)}_{\Cell{#2}{#3}}}
\newcommand{\Vdr}[3]{\cV^{(#1)}_{\Cell{#2}{#3}}}
\newcommand{\vVdr}[3]{\vec{\cV}^{(#1)}_{\Cell{#2}{#3}}}
\newcommand{\F}[2]{F_{\Cell{#1}{#2}}}
\newcommand{\z}[2]{\zeta_{\Cell{#1}{#2}}}
\newcommand{\g}[2]{g_{CNO_{\Cell{#1}{#2}}}}
\newcommand{\N}[1]{\cN^{(#1)}}
\newcommand{\R}[2]{R_{\Cell{#1}{#2}}}
\newcommand{\Co}[4]{\Gamma^{\Cell{#1}{#2}}_{\Cell{#3}{#4}}}
\newcommand{\pr}[4]{\Pi^{\Cell{#1}{#2}}_{\Cell{#3}{#4}}}
\newcommand{\W}[4]{W^{\Cell{#1}{#2}}_{\Cell{#3}{#4}}}
\newcommand{\Wp}[4]{W'^{\Cell{#1}{#2}}_{\Cell{#3}{#4}}}
\newcommand{\Di}[2]{D_{#1}^{(#2)}}
\newcommand{\Dip}[2]{D_{#1}^{'(#2)}}
\newcommand{\Gp}[4]{G^{\Cell{#1}{#2}}_{\Cell{#3}{#4}}}
\newcommand{\gp}[4]{g^{\Cell{#1}{#2}}_{\Cell{#3}{#4}}}
\newcommand{\alp}[2]{\alpha^{#1}_{#2}}
\renewcommand{\t}[2]{\tau^{#1}_{#2}}
\newcommand{\Ps}[4]{P^{\Cell{#1}{#2}}_{\Cell{#3}{#4}}}
\newcommand{\Amp}[2]{w^{#1}_{#2}}
\newcommand{\sig}[2]{\sigma^{#1}_{#2}}
\newcommand{\Ist}[2]{I^{(c)}_{\Cell{#1}{#2}}}
\newcommand{\Pst}[2]{P^{(c)}_{\Cell{#1}{#2}}}
\newcommand{\Gap}[2]{P^{(g)}_{\Cell{#1}{#2}}}
\newcommand{\IG}[2]{I^{(g)}_{\Cell{#1}{#2}}}
\newcommand{\D}[4]{d\bra{\Cell{#1}{#2}, \, \Cell{#3}{#4}}}
\newcommand{\Dd}[4]{d^2\bra{\Cell{#1}{#2}, \, \Cell{#3}{#4}}}
\newcommand{\cA}{{\mathcal A}}
\newcommand{\cB}{{\mathcal B}}
\newcommand{\cC}{{\mathcal C}}
\newcommand{\cD}{{\mathcal D}}
\newcommand{\cE}{{\mathcal E}}
\newcommand{\hcE}{{\hat{\mathcal E}}}
\newcommand{\cG}{{\mathcal G}}
\newcommand{\cK}{{\mathcal K}}
\newcommand{\cH}{{\mathcal H}}
\newcommand{\cI}{{\mathcal I}}
\newcommand{\cL}{{\mathcal L}}
\newcommand{\cLT}{{\mathcal L}^T}
\newcommand{\tcL}{\tilde{\cL}}
\newcommand{\cM}{{\mathcal M}}
\newcommand{\cN}{{\mathcal N}}
\newcommand{\cP}{{\mathcal P}}
\newcommand{\hcP}{\hat{\mathcal P}}
\newcommand{\pP}{{\mathbb P}}
\newcommand{\cQ}{{\mathcal Q}}
\newcommand{\cO}{{\mathcal O}}
\newcommand{\tcP}{\tilde{\cP}}
\newcommand{\cR}{{\mathcal R}}
\newcommand{\cS}{{\mathcal S}}
\newcommand{\cF}{{\mathcal F}}
\newcommand{\cX}{{\mathcal X}}
\newcommand{\cU}{{\mathcal U}}
\newcommand{\cV}{{\mathcal V}}
\newcommand{\cT}{{\mathcal T}}
\newcommand{\cTm}{{\mathcal T}^{-1}}
\newcommand{\cW}{{\mathcal W}}
\newcommand{\fA}{\mathfrak{A}}
\newcommand{\fB}{\mathfrak{B}}
\newcommand{\fG}{\mathfrak{G}}
\newcommand{\vzA}{\vec{\zeta}_A}
%\newcommand{\veta}{\vec{\eta}}
\newcommand{\veta}{\boldsymbol{\eta}}
\newcommand{\vz}{\vec{0}}
\newcommand{\vU}{\bold{U}}
\newcommand{\vun}{\vec{1}_N}
\newcommand{\vzG}{\vec{\zeta}_G}
%\newcommand{\vxi}{\vec{\xi}}
\newcommand{\veps}{\boldsymbol{\epsilon}}
\newcommand{\vxi}{\boldsymbol{\xi}}
\newcommand{\vxiT}{\boldsymbol{\xi}^T}
\newcommand{\vth}{\boldsymbol{\theta}}
\newcommand{\vtha}{\boldsymbol{\theta}^\ast}
\newcommand{\vcP}{\vec{\cP}}
\newcommand{\vpP}{\vec{\pP}}
\newcommand{\vcQ}{\vec{\cQ}}
\newcommand{\vm}{\vec{m}}
\newcommand{\vk}{\vec{k}}
\newcommand{\vV}{\vec{V}}
\newcommand{\vVA}{\vV_\fA}
\newcommand{\vVG}{\vV_\fG}
\newcommand{\VAa}{V_{\fA,a}}
\newcommand{\vVB}{\vV_\fB}
\newcommand{\VBa}{V_{\fB,\alpha}}
\newcommand{\vA}{\vec{A}}
\newcommand{\ket}[1]{|#1 \rangle}
\newcommand{\Ket}[2]{|#1_#2 \rangle}
\newcommand{\Bra}[2]{\langle #1_#2|}
\newcommand{\Proj}[2]{|#1_#2\rangle \langle #1_#2|}
\newcommand{\Scal}[4]{\langle #1_#2|#3_#4\rangle}
\newcommand{\vphi}{\vec{\phi}}
\newcommand{\kphi}[1]{\Ket{\phi}{#1}}
\newcommand{\bphi}[1]{\Bra{\phi}{#1}}
\newcommand{\vpsi}{\vec{\psi}}
\newcommand{\vPhi}{\vec{\Phi}}
\newcommand{\kpsi}[1]{\Ket{\psi}{#1}}
\newcommand{\bpsi}[1]{\Bra{\psi}{#1}}
\newcommand{\Hh}[2]{\cH^{#1}_{#2}}

\newcommand\bloc[2]{{\omega}_{#1}^{#2}}
\newcommand\blocp[2]{{\omega'}_{#1}^{#2}}
\newcommand\sif[1]{\seq{\omega}{-\infty}{#1}}
\newcommand{\Pnc}[2]{\mathds{P}_n\bra{#1 \, \left| \, #2 \right.}}
\newcommand\seq[3]{{#1}_{#2}^{#3}}
\newcommand{\probc}{\mathds{P}}
\newcommand{\Probc}[2]{\mathds{P}\bra{#1 \, \left| \, #2 \right.}}
\newcommand\moy[1]{\mu\bra{#1}}
\newcommand\moysp[1]{\mu^{(sp)}\bra{#1}}
\newcommand\musp{\mu^{(sp)}}
\newcommand{\ent}[1]{\left[\, #1 \, \right]}
\newcommand\gk[1]{g_k\pare{#1}}
\newcommand\Csp[2]{\mathcal{C}^{(sp)}\bra{#1, #2}}
\newcommand\sk[1]{\sigma_k({#1})}
\newcommand\sj[1]{\sigma_j({#1})}
\newcommand\nkrm{\tau_k(r-1,.)}

\newcommand{\re}[1]{Re\pare{#1}}
\newcommand{\im}[1]{Im\pare{#1}}

%__________________________________________
% newcommands Luc
%__________________________________________

\newcommand{\bea}{\begin{eqnarray*}}
\newcommand{\eea}{\end{eqnarray*}}
\newcommand{\beq}{\begin{eqnarray}}
\newcommand{\eeq}{\end{eqnarray}}



\def\rad{\stackrel{\rm d}{\ra}}

%\input{defpera.tex}
\def\dd{\mathrm{d}}
\def\arg{\mathrm{arg}}
\def\Arg{\mathrm{Arg}}
\def\dim{\mathrm{dim}}
\def\diag{\mathrm{diag}}
\def\diam{\mathrm{diam}}
\def\Pr{\mathrm{Prob}}
\def\mp{\partial}
\def\rank{\mathrm{rank}}
\def\tr{\mathrm{trace}}
\def\vol{\mathrm{vol}}

\def\mt{\theta}

\def\Ex{\mathsf{E}}
\def\Var{\mathsf{Var}}
\def\var{\mathsf{var}}

\def\tDelta{\widetilde{\Delta}}
\def\bDelta{\overline{\Delta}}

\def\ma{\alpha}
\def\mg{\gamma}
\def\ml{\lambda}
\def\mL{\Lambda}
\def\mt{\theta}
\def\mth{\hat\theta}
\def\me{\epsilon}
\def\mve{\varepsilon}
\def\mdb{\boldsymbol{\delta}}
\def\ms{\sigma}
%\def\Ex{{\rm I}\!{\rm E}}
%\def\Ex{\mathbb{E}}

\def\ra{\rightarrow}
\def\TT{^\top}
\def\e1{\mathsf{e}}

\def\Ab{\mathbf{A}}
\def\Db{\mathbf{D}}
\def\bDb{\overline{\mathbf{D}}}
\def\Eb{\mathbf{E}}
\def\Fb{\mathbf{F}}
\def\Gb{\mathbf{G}}
\def\Hb{\mathbf{H}}
\def\Ib{\mathbf{I}}
\def\Jb{\mathbf{J}}
\def\Kb{\mathbf{K}}
\def\bK{\overline{K}}
\def\bKb{\overline{\mathbf{K}}}
\def\bmOb{\overline{\boldsymbol{\Omega}}}
\def\bmGb{\overline{\boldsymbol{\Gamma}}}
\def\mob{\boldsymbol{\omega}}
\def\bmob{\overline{\boldsymbol{\omega}}}
\def\tK{\widetilde{K}}
\def\tKb{\widetilde{\mathbf{K}}}
\def\Mb{\mathbf{M}}
\def\bMb{\overline{\mathbf{M}}}
\def\Ob{\mathbf{O}}
\def\Qb{\mathbf{Q}}
\def\Rb{\mathbf{R}}
\def\Sb{\mathbf{S}}
\def\Tb{\mathbf{T}}
\def\Xb{\mathbf{X}}
\def\Vb{\mathbf{V}}
\def\Zb{\mathbf{Z}}
\def\Wb{\mathbf{W}}
\def\mLb{\boldsymbol{\Lambda}}
\def\mOb{\boldsymbol{\Omega}}
\def\Phib{\boldsymbol{\Phi}}
\def\Sigmab{\boldsymbol{\Sigma}}

\def\1b{\mathbf{1}}
\def\0b{\mathbf{0}}
\def\bb{\mathbf{b}}
\def\cb{\mathbf{c}}
\def\Cb{\mathbf{C}}
\def\db{\mathbf{d}}
\def\eb{\mathbf{e}}
\def\fb{\mathbf{f}}
\def\gb{\mathbf{g}}
\def\hb{\mathbf{h}}
\def\kb{\mathbf{k}}
\def\pb{\mathbf{p}}
\def\qb{\mathbf{q}}
\def\rb{\mathbf{r}}
\def\tb{\mathbf{t}}
\def\ub{\mathbf{u}}
\def\vb{\mathbf{v}}
\def\wb{\mathbf{w}}
\def\xb{\mathbf{x}}
\def\yb{\mathbf{y}}
\def\Yb{\mathbf{Y}}
\def\zb{\mathbf{z}}
\def\alphab{\boldsymbol{\alpha}}
\def\betab{\boldsymbol{\beta}}
\def\gammab{\boldsymbol{\gamma}}
\def\etab{\boldsymbol{\eta}}
\def\phib{\boldsymbol{\phi}}
\def\taub{\boldsymbol{\tau}}
\def\mveb{\boldsymbol{\varepsilon}}
\def\mtb{\boldsymbol{\theta}}
\def\zetab{\boldsymbol{\zeta}}

\def\SA{{\mathscr A}}
\def\SB{{\mathscr B}}
\def\SE{{\mathscr E}}
\def\SF{\mathcal{F}}
\def\SSF{\mathscr{F}}
\def\SH{\mathcal{H}}
\def\SSL{\mathds{L}}
\def\SM{{\mathscr M}}
\def\SN{{\mathscr N}}
\def\SO{\mathcal{O}}
\def\SP{{\mathscr P}}
\def\SU{{\mathscr U}}
\def\SS{{\mathscr S}}
\def\ST{{\mathscr T}}
\def\SX{{\mathscr X}}




\newcommand{\fin}{\mbox{}~\hfill\mbox{$\lhd$}}
\newcommand{\carre}{\mbox{}~\hfill\rule{2mm}{2mm}}


\newtheorem{rmk}{Remark}
\newtheorem{defi}{Definition}
\newtheorem{thm}{Theorem}
\newtheorem{lemma}{Lemma}
\newtheorem{corollary}{Corollary} %Ludovic
\newtheorem{example}{Example}
%\newdefinition{lemma}{Lemma}
%\newtheorem{propo}{Proposition}

\newcommand{\vsp}{\vspace{0.3cm}}
\newcommand{\lline}{\rule{10cm}{0.2mm}}







%\addbibresource{bibliography.bib} % The filename of the bibliography
\title{Parameters estimation in a retinal network model}


\author{
 Bruno Cessac, Ludovic Sacchelli, Luc Pronzato
}


\begin{document}

\setlength\parindent{24pt}

\maketitle

\section{Context}

We consider a linear, stochastic, dynamical system:
%
\begin{equation} \label{eq:SD}
d\vX\,=\,\cL.\vX \,dt \,+\,\vF(t) \,dt \,+\, \sigma_B \,d\vW(t),
\end{equation}
%
where $\vX \in \setR^N$ is the state vector, with entries $X_i$, $i=1 \dots N$, $\cL$ is a $N \times N$ matrix. We assume $\cL$ to be diagonalisable, with eigenvalues
$\lambda_n, n=1 \dots N$, left (normalised) eigenvectors $\tvw_n$ - where $T$ denotes the transpose - and right normalised eigenvectors $\vv_n$.
We note $\Lambda$ the diagonal form of $\cL$ so that $\Lambda\,=\,\cP^{-1}.\cL.\cP$ where $\cP^{-1}$ is the matrix whose rows are the
left eigenvectors $\tvw_n$ while $\cP$ is the matrix whose columns are the right eigenvectors $\vv_n$. 
We have thus (spectral decomposition theorem):
%
\begin{equation}\label{eq:SpectralDecomposition}
\cL\,=\,\sum_{n=1}^N \lambda_n \, \vv_n.\tvw_n,
\end{equation} 
%
and:
%
\begin{equation}\label{eq:Orthogonality}
\tvw_m.\vv_n\,= \, \delta_{mn},
\end{equation}
%
where $\delta_{mn}$ is the Kronecker symbol. Note that the $\vv_n$s are not orthogonal in general, neither nor the $\tvw_n$s. Moreover, eigenvalues and eigenvectors are in general complex. We assume that all $\cL$ is stable, that is, its eigenvalues have a strictly negative real part. More precisely, we order the eigenvalues such that:
%
\begin{equation}\label{eq:OrderEigenvalues}
\re{\lambda_N} \leq \re{\lambda_{N-1}} \leq \dots \leq \re{\lambda_1} < \gamma,
\end{equation}
%
for some $\gamma < 0$. As a consequence, $\cL$ is invertible.

In eq. \eqref{eq:SD}, $\vF(t)$ is a smooth time dependent vector, called from now on the "input", with entries $F_i(t)$, $i=1 \dots N$. $d\vW$ is a Brownian noise, whose entries $dW_{i}(t)$ obey $\Exp{\vW}{dW_{i}(t)}=0$,
$\Exp{\vW}{dW_{i}(t)dW_{i'}(t')}\,=\,\delta(t-t')\delta_{ii'}$ where we note
$\Exp{\vW}{}$ the expectation with respect to $\vW$.\\

We assume that the dynamical system \eqref{eq:SD} was initially prepared in a distant past, $t_0$, and that the input
$\vF$ is applied starting from $t=0$. The solution of \eqref{eq:SD} is then:
%
\begin{equation}\label{eq:SolSDGen}
\vX(t)\,=\, e^{\cL(t-t_0)}.\vX(t_0) \,+\,  \int_{0}^t e^{\cL(t-s)}.\vF(s) ds \,+\, \sigma_B \, \int_{t_0}^t e^{\cL(t-s)}.d\vW(s). 
\end{equation}
%
Using the stability of $\cL$ we let the initial time $t_0$ tend to $-\infty$ so as to remove the dependence on the initial condition, $\vX(t_0) $. The solution becomes, for $t>0$:
%
\begin{equation}\label{eq:SolSD}
\vX(t)\,=\, \int_{0}^t e^{\cL(t-s)}.\vF(s) ds \,+\, \sigma_B \, \int_{-\infty}^t e^{\cL(t-s)}.d\vW(s). 
\end{equation}
%
This is a Gaussian stochastic process with mean:
%
\begin{equation}\label{eq:Exp_X}
\mu(t) \deq \Exp{\vW}{\vX(t)} \,=\, \int_{0}^t e^{\cL(t-s)}.\vF(s) ds, 
\end{equation}
%
and covariance:
%
\begin{equation}\label{eq:Cov_X}
\cC(t,t') \deq \Exp{\vW}{\pare{\vX(t)-\mu(t)}.\pare{\vX(t')-\mu(t')}^T} \,=\, \sigma_B^2 \, \int_{-\infty}^t e^{\cL(t-s)}.e^{\cLT(t'-s)} \, ds, 
\end{equation}
%
for $t\leq t'$. Note that $\cL$ does in general  not commute with its transpose. The covariance can be expressed in terms of eigenvalues and eigenvectors using the spectral decomposition theorem \eqref{eq:SpectralDecomposition}:
%
%$$
%\cC(t,t') \,=\, \int_{-\infty}^t \pare{\sum_{n=1}^N e^{\lambda_n(t-s)} \, \vv_n.\tvw_n}.\pare{\sum_{m=1}^N e^{\lambda_m(t-s)} \, \vw_m.\tvv_m} \,ds \,=\, \sum_{n,m=1}^N \,\vv_n.\tvw_n.\vw_m.\tvv_m \int_{-\infty}^t e^{\pare{\lambda_n+\lambda_m}(t-s)} \, ds
%$$
%
\begin{equation}\label{eq:Cov_X_Spectral}
\cC(t,t') \,=\, -\, \sum_{n=1}^N \,\vv_n.\tvw_n. \sum_{m=1}^N \frac{e^{\lambda_m(t'-t)}}{\lambda_n+\lambda_m}.\vw_m.\tvv_m
\,=\, -\,  \sum_{m=1}^N e^{\lambda_m(t'-t)} \sum_{n=1}^N \, \frac{\vv_n.\tvw_n.\vw_m.\tvv_m}{\lambda_n+\lambda_m}
\end{equation}
%%
Remember that the eigenvalues have a negative real part, and, when complex, go by conjugate pairs. So, correlation is real, as expected. It depends only on the difference $t'-t$ (stationarity) because we took the limit of the initial time tending to $-\infty$. 

The derivation of the covariance is done in the appendix \ref{Sec:ACov}.
A simpler form can be obtained introducing the metric:
%
\begin{equation}\label{eq:Metric}
\cG\,=\,\cP^{-1}.\pare{\cP^{-1}}^T,
\end{equation}
%
with entries:
\begin{equation}\label{eq:gij}
g_{mn}\,=\,<\vw_m,\vw_n>,
\end{equation}
%
and the matrix:
%
\begin{equation}\label{eq:Delta}
\cD(t,t')\,=\,\cP^{-1}.\cC(t,t').\pare{\cP^{-1}}^T.
\end{equation}

Then:
%
\begin{equation}\label{eq:Deltamn}
\cD_{mn}(t,t') \,=\,  
- \, \sigma_B^2 \, \frac{g_{mn}}{\lambda_m \,+\, \lambda_n} \, e^{\lambda_n \, \pare{t'-t}}.
\end{equation}

Defining the matrix $\cA$ with entries:
%
\begin{equation}\label{eq:Amn}
A_{mn}=\frac{g_{mn}}{\lambda_m \,+\, \lambda_n},
\end{equation}
%
we have:
%
\begin{equation}\label{eq:D_mat}
\cD(t,t') \,=\,- \, \sigma_B^2 \cA.e^{\Lambda(t-t')}.
\end{equation}

It is also useful to compute the inverse of $\cD(t,t')$, $\cD^{-1}(t,t')\,=\ \cP^T.\cC^{-1}(t,t').\cP$, which is:
%
\begin{equation}\label{eq:Dmoins1}
\cD^{-1}(t,t') \,=\, - \, \sigma_B^{-2} \, e^{-\Lambda(t-t')}.\cA^{-1},
\end{equation} 
%
with entries:
%
\begin{equation}\label{eq:Dmoins1}
\cD^{-1}_{mn}(t,t') \,=\, - \, \sigma_B^{-2} \, e^{-\lambda_m(t-t')}.\cA^{-1}_{mn},
\end{equation} 

Note that this matrix diverges as $t-t' \to +\infty$, that is, the longer the time interval $t-t'$, the larger the weight of $\cD^{-1}(t,t')$.

\Luc{L'article assez récent \cite{duarte2021optimal} contient pas mal de références pour le problème de choix d'une excitation optimale  pour estimer les paramètres d'un modèle dynamique. Ce n'est pas forcément le plus clair sur le sujet, mais on devrait peut être indiquer en quoi le problème que l'on considère est différent du leur.}

% \B{Avec cette matrice, le critère $\phi$ fait intervenir la forme
% quadratique $\Exp{\xi}{\sum_{l,l',m,n} \eta_m(l) \, \cD^{-1}_{mn}(t,t')\,\eta_n(l')} $
% $=- \, \sigma_B^{-2} \, \sum_{l,l',m,n} e^{-\lambda_m(t-t')}\,\cA^{-1}_{mn} \, \Exp{\xi}{Z_m(l) \,Z_n(l')}$
% où $\vZ(l)\equiv \zeta(\vth,\vu)-\zeta(\vth^0,\vu)$.
% On fait donc intervenir la correlation, $\cC_{\vZ}(l,l')$, de l'erreur $\vZ$ sous
% $\xi$, soit les matrice d'entrées $\Exp{\xi}{Z_m(l) \,Z_n(l')}$.
% Le critère s'écrit alors (le facteur $- \, \sigma_B^{-2}$ n'intervenant pas dans la minimisation):
% %
% $$
% \phi(\xi)= \min_{\vth} \frac{\sum_{l,l',m,n} e^{-\lambda_m(t-t')}\,\cA^{-1}_{mn} \, \Exp{\xi}{Z_m(l) \,Z_n(l')}}{\|\vth \,-\, \vth^0\|^2},
% $$
% %
% où tous les termes du dénominateur dépendent de $\vth$.
% Cela fait apparaitre (1) la métrique $G$ induites par les vecteurs propres; (2) les valeurs propres; (3) les correlations de l'erreur sous $\xi$.
% }


%\TBC
\section{Estimation problem}

The matrix $\cL$ depends on a set of parameters, noted in the form of a vector $\bar\mtb$ of dimension $P$. The problem is to estimate at best these parameters by using a loop strategy: (1), Present a series of inputs $\vF^{(k,1)}(t), \, k=1 \dots K$, distributed according to a probability density $\xi^{(1)}$;
(2), Estimate $\bar\mtb$ by minimizing some distance between the measured response and the theoretical response \eqref{eq:SolSD}; (3), From this estimation propose a new series of input $\vF^{(k,2)}(t), \, k=1 \dots K$, optimal for some criterion, and distributed according to a probability density $\xi^{(2)}$. And so on. 
In the following, we use the methodology proposed in \cite{pronzato:08} and adapt our case to this setting.

\subsection{Input structure and notations}\label{eq:InputStructure}

Although time is continuous, observations are made on a discrete time with a step $\delta>0$, over a time horizon $[0,(H-1)\delta]$. Thus, the input $\vF$ is a piecewise constant function\footnote{We drop here the dependence in $k$ to alleviate the notations.} of time taking the value $\vF(t_l)$ on the time interval $[t_l, t_{l+1}[$, with $t_l=(l-1)\delta$. $\vF$ can therefore be represented as a list of $H$ vectors of dimension $N$, $\left\{\vF(t_0), \dots, \vF(t_{H-1})\right\}$ that we convert into a  one dimensional vector, $\vu$, of dimension $Q=N \times H$, with entries $u_q\,=\,F_i(t_l)$, with $q=i+(l-1)N$. 

We assume here, for simplicity, that $\delta$ is quite smaller than the characteristic times in the problem so that we may approximate the differential equation \eqref{eq:SD} by an Euler discretisation:
%
$$
\vX(t_l) \,-\, \vX(t_{l-1})\,=\, \delta \cL \vX(t_{l-1}) \,+\, \delta \vF(t_{l-1}) \,+\, \sigma_B \vN(t_{l}),
$$
%
where $\vN(t_{l})$ is a Gaussian vector, centered, with covariance
%
$$
\Exp{}{\vN(t_{l}).\vN^T(t_{l'})} \,=\, \delta \cI_N \delta_{l l'},
$$
%
$\cI_N$ being the $N$ dimensional identity matrix and $\delta_{ll'}$ the Kronecker symbol.
% %
Without loss of generality, because the system is linear, we can assume that the stimulation starts with the initial condition $\vX(0)=0$. In addition, like in the previous section, we assume that the noise was applied in a distant past $\delta n_0$ where $n_0 \to -\infty$. Then:
%
\begin{equation}\label{eq:SDD}
\vX(t_l)  \,=\, \delta  \sum_{n=0}^{l-1} \cM^{l-n-1} \vF(t_n) \,+\, \sigma_B \sum_{n=n_0}^{l}  \cM^{l-n} \vN(t_{n})
\end{equation}
%
where $\cM \,=\, \cI_N \,+\, \delta \cL$.
%
With a slight abuse of notation\footnote{Because $\vX(t_l)$ does not depend on the whole input vector $\vu$, only on the entries corresponding to times $< t_l$} we note $\yb(t_\ell,\ub) \,=\,\vX(t_l)$. We also note $\etab(t_\ell,\bar\mtb,\ub) \,=\, \delta  \sum_{n=0}^{l-1} \cM^{l-n-1} \vF(t_n)$ and $\mveb(t_\ell) \,=\, \sigma_B \sum_{n=n_0}^{l}  \cM^{l-n} \, \vN(t_{n})$ so that \eqref{eq:SDD} reads:
%
\begin{equation}\label{eq:SDE}
\yb(t_\ell,\ub) = \etab(t_\ell,\bar\mtb,\ub)+\mveb(t_\ell),
\end{equation}
%
where $\mveb(t_\ell)$ is called the "error".

 As a sum of Gaussian centered variables the error $\mveb(t_\ell)$ is Gaussian, centered. Its covariance $\Sigma(t_l,t_{l'}) \,=\, \Exp{}{\mveb(t_\ell)\mveb^T(t_{l'})}$ obeys:
%
$$
\Sigma(t_l,t_{l'}) \,=\, \sigma_B^2 \, \sum_{n=n_0}^{l} \sum_{n'=n_0}^{l'} \cM^{l-n} . \underbrace{\Exp{}{\vN(t_{n}).\vN^T(t_{n'})}}_{\delta \cI_N \delta_{n n'}} . \pare{\cM^T}^{l'-n'}
\,=\, \delta \,\sigma_B^2 \,  \sum_{n=n_0}^{l} \cM^{l-n} . \pare{\cM^T}^{l'-n},
$$
%
assuming $l \leq l'$. Using the diagonal form of $\cM$ this is:
%
$$
\Sigma(t_l,t_{l'}) \,=\, \delta \,\sigma_B^2 \,  \cP \pare{\sum_{n=n_0}^{l} \pare{\cI_N\,+\, \delta \Lambda}^{l-n}.\underbrace{\cP^{-1}.\cP^{{-1}^T}}_{\cG}.\pare{\cI_N\,+\, \delta \Lambda}^{l'-n}}.\cP^T
$$
%
Noting $\cD(t_l,t_{l'}) \,=\,\cP^{-1}.\Sigma(t_l,t_{l'}).{\cP^T}^{-1}$ and denoting $D_{ik}(t_l,t_{l'})$ the entry $i,k$ of the $N \times N$ matrix $\cD(t_l,t_{l'})$, we have:
%
$$
D_{ik}(t_l,t_{l'}) \,=\, \delta \,\sigma_B^2 \,  g_{ik} \, \sum_{n=n_0}^{l} \pare{1+\delta \lambda_i}^{l-n}\pare{1+\delta \lambda_k}^{l'-n}
$$
%
$$
\,=\, \delta \,\sigma_B^2 \,  g_{ik} \, \pare{1+\delta \lambda_i}^{l}\pare{1+\delta \lambda_k}^{l'} \, \sum_{n=n_0}^{l} \pare{1+\delta \lambda_i}^{-n}\pare{1+\delta \lambda_k}^{-n}.
$$
%
We set $q=\pare{1+\delta \lambda_i}\pare{1+\delta \lambda_k} = 1 + \delta \pare{\lambda_i+\lambda_k} + \delta^2 \lambda_i\lambda_k$. We assume that $\delta$ is small enough so that $q \sim 1 + \delta \pare{\lambda_i+\lambda_k}$ and that $|q|<1$ (which holds for sufficiently small $\delta$ as all eigenvalues are stable\footnote{
$\abs{1+\delta\lambda}^2 = \abs{1+\delta \lambda_r + i \delta \lambda_i}^2$ 
$=\pare{1+\delta \lambda_r}^2 + \delta^2 \lambda_i^2  = 1 +2 \delta \lambda_r + \delta^2\pare{\lambda_r^2+\lambda_i^2}  < 1$  
$\Leftarrow \delta < - \frac{2 \lambda_r}{\lambda_r^2+\lambda_i^2}$ }
).
%
Then,
%
$$\sum_{n=n_0}^{l} \pare{1+\delta \lambda_i}^{-n}\pare{1+\delta \lambda_k}^{-n}
=\sum_{n=n_0}^{l} q^{-n} 
=q^{-n_0} \frac{1-q^{n_0-l+1}}{1-q}
=\frac{q^{-n_0}-q^{-l+1}}{1-q} 
 \to
- \frac{q^{-l+1}}{1-q}
,$$
%
as $n_0 \to -\infty$.\\
%
It follows that:
%
$$
D_{ik}(t_l,t_{l'}) \,=\, \delta \,\sigma_B^2 \,  g_{ik} \,
\pare{1+\delta \lambda_i}^{l}\pare{1+\delta \lambda_k}^{l'} \, \frac{\pare{1+\delta \lambda_i}^{-l+1}\pare{1+\delta \lambda_k}^{-l+1}}{\delta \pare{\lambda_i+\lambda_k} + \delta^2 \lambda_i\lambda_k}
$$
%
$$
\,=\, \sigma_B^2 \,  g_{ik} \, \frac{\pare{1+\delta \lambda_i}\pare{1+\delta \lambda_k}}{\pare{\lambda_i+\lambda_k} + \delta \lambda_i\lambda_k} \, \pare{1+\delta \lambda_k}^{l'-l},
$$
%
which is similar to \eqref{eq:Deltamn}.
%
% $$
% \cD_{mn}(t,t') \,=\,  
% - \, \sigma_B^2 \, \frac{g_{mn}}{\lambda_m \,+\, \lambda_n} \, e^{\lambda_n \, \pare{t'-t}}.
% $$
%
Note that the covariance is therefore not diagonal, although its off diagonal terms decay exponentially fast.

% The time discretized response to the input $\vF$ reads then, from \eqref{eq:SolSDGen}, for $l \in \Set{0, \dots, H-1}$:
% %
% $$
% \vX(t_l)\,=\, 
% %e^{\cL \delta}.\vX(l\delta) \,+\,  
% \int_{0}^{t_l} e^{\cL(t_l-s)}.\vF(s) ds \,+\,\sigma_B \, \int_{-\infty}^{t_l} e^{\cL(t_l-s)}.d\vW(s). 
% $$
% %

% Using the fact that $\vF$ is piecewise constant this gives:
% % 
% \begin{equation}\label{eq:DiscX}
% \vX(l\delta)\,=\, \,\underbrace{\sum_{r=0}^{l-1} \,\cI^{l,r}.\vF(r\delta)}_{\vD(l\delta)} \,+\, \underbrace{\sigma_B \, \int_{0}^{l\delta} e^{\cL(l\delta-s)}.d\vW(s)}_{\vN(l\delta)},
% \end{equation}
% %
% where:
% %
% \begin{equation}\label{eq:Ilr}
% \cI^{l,r}\,=\,\int_{\delta r}^{\delta (r+1)} e^{\cL\pare{l \delta - s}} \,ds
% \end{equation}
% %
% is a $N \times N$ matrix with entries $\cI^{l,r}_{ij}$. We have separated \eqref{eq:DiscX} into a deterministic term, $\vD(l\delta)$,  and a noisy term, $\vN(l\delta)$ which does not depend on the input index. 
% The deterministic response
% is thus a list of vectors $\left\{\vD(\delta), \dots, \vD(H\delta)\right\}$ that we convert  into a $Q$-dimensional vector  $\veta(\vth,\vu)$: 
% That is:
% %
% \begin{equation}\label{eq:eta}
% \eta_q(\vth,\vu)\,=\,D_i(l\,\delta), \, l=1 \dots H.
% \end{equation}
% %
% Compared to the definition of $\vu$ there is a shift in time: the index $q$ of $\vu$ corresponds to $q=i+(l-1)N$ while
% the index $q$ of $\veta$ corresponds to $q=i+lN$. This is an artifact coming from the time discretisation. This is transparent, however,
% when using the vector forms.

% The noise vector, $\vN(l\delta)$, is Gaussian, centered, with a correlation given by:
% %
% \begin{equation}\label{eq:Cov_eps}
% \cC(l,l')\,=\, -\,  \sigma_B^2 \, \sum_{m=1}^N e^{\lambda_m(l'-l)\delta} \sum_{n=1}^N \, \vv_n.\tvw_n.\vw_m.\tvv_m \frac{1\,-\,e^{(\lambda_n+\lambda_m)\delta}}{\lambda_n+\lambda_m},
% \end{equation}
% %
% for $l\leq l'$ using the same method as for \eqref{eq:Cov_X}. Note that this "noise" contains correlations informative on the network structure as it depends on the eigen-elements of $\cL$, but it does not depend on the input.

% In the same way as for $\veta$ above, we order the noise vectors at successive times in a $Q$-dimensional vector $\veps$ with
% entries $\epsilon_q\,=\,N_i(l\delta)$. The $\epsilon_q$s are not independent, as they are correlated through dynamics, via the correlation \eqref{eq:Cov_eps}. 
% It will be nevertheless useful, for estimation, to assume that they are independent and Gaussian. This holds if the correlation decay is sufficiently fast,
% that is, using \eqref{eq:Cov_eps}, if $\gamma\delta \ll 1$, where $\gamma$ is defined in \eqref{eq:OrderEigenvalues}. In this case, we note $\sigma^2$ the variance of $\epsilon_q$s (given by \eqref{eq:Cov_eps} with $l=l'$),
% and $\phi$ their marginal probability density. 

% The vector $\eta_q(\vth,\vu)$ is what one would observe if the system had the evolution \eqref{eq:SD} with a matrix $\cL$ having the parameters set $\vth$. The actual response is a vector $\vy$, of dimension $Q$, in general different from $\eta_q(\vth,\vu)$. We however assume that $\vy$ has been produced by a system of the form \eqref{eq:SD}, where $\cL$ has a set of parameters $\vtha$ that we seek to determine. Thus, for each input $k$:
% %
% \begin{equation}\label{eq:y}
% \vy \,=\, \eta_q(\vtha,\vu) \,+\, \veps.
% \end{equation}

% \subsection{Estimation}\label{Sec:Estimation}

% We assume now that the inputs are sampled from a probability distribution with density $\xi(\vu)$.
% We are given $K$ samples of the input, corresponding to the $Q$-dimensional vectors $\vu^{(k)}$, $k=1\dots K$.

% \subsubsection{Criteria}\label{Sec:Criteria}
% The estimation problem consists of finding (least square estimator):
% %
% \begin{equation}\label{eq:theta_hat}
% \hat{\vth}_{WLS} \,=\, \argmin_{\vth} \, \frac{1}{K} \, \sum_{k=1}^K w(\vu^{(k)}) \|\vy \,-\, \eta_q(\vth,\vu^{(k)}) \|^2,
% \end{equation}
% %
% where $w$ is a known, bounded, function. Conditions for the almost-sure convergence (under $\xi$) of this estimator as the number of samples $K$ tends to $+\infty$
% are stated in \cite{pronzato:08}. Under these conditions $\hat{\vth}_{WLS}$ also obeys a central limit theorem whose variance is minimal 
% when $w()\,=\,c\sigma^{-2}$.

% When the $\epsilon_q$s are i.i.d. another possible estimator is the maximum-likelihood. When the $\epsilon_q$s are Gaussian, the two estimators coincide. 

% \subsubsection{The Fisher Matrix}\label{Sec:Fisher}

% The estimator obeys a central limit theorem whose covariance is the Fisher matrix \BL{Est-ce correct en dimension $N$ ?}
% %
% \begin{equation}\label{eq:Fisher}
% \cM_F(\xi,\vth) = \int_{\cU} \pare{\frac{\partial \veta(\vth,\vu)}{\partial \vth}}^T.\pare{\frac{\partial \veta(\vth,\vu)}{\partial \vth}} \,\xi(d\vu).
% \end{equation}

% This matrix will be also used below for the input design.
% It is computed in the appendix. Its entries take the general form:
% %
% %
% \begin{equation}\label{eq:MF_exp}
% \cM_{pp'}(\xi,\vth) \,=\,  \sum_{l=1}^H \sum_{r,r'=0}^{l-1} \, \sum_{i=1}^N \,\sum_{j,j'=1}^N \, \frac{\partial \cI^{l,r}_{ij}}{\partial \theta_p}\,\Delta_{jj'}(r,r')\, \frac{\partial \cI^{l,r'}_{ij'}}{\partial \theta_{p'}}\, 
% \end{equation}
% %
% where
% %
% \begin{equation}\label{eq:Delta}
% \Delta_{jj'}(r,r') \,=\, \Exp{\xi}{F_j(r\delta)\,F_{j'}(r'\delta) }.
% \end{equation}


% %
% \paragraph{Remarks}

% \begin{itemize}

% \item In the Fisher matrix computation  there is a contribution,
% \eqref{eq:Delta}, which depends only on the input distribution, and  contributions, $\frac{\partial \cI^{l,\bullet}_{ij}}{\partial \theta_\bullet}$, which depends only on $\cL$ and thereby, on $\vth$ (see eq. \eqref{eq:dIdtheta} in the appendix). %This separation, is the key toward the implementation of the method sketched in the beginning of this section.
% \item Note that, in the series \eqref{eq:dIdtheta} defining $\frac{\partial \cI^{l,\bullet}_{ij}}{\partial \theta_\bullet}$ the convergence is controlled by $\delta\Lambda$
% that is by $\delta\,\gamma$ which as to be $\ll 1$ so as to be numerically tractable. However, $\delta$ is constrained by the experimental set up, while $\gamma$, the upper bound on the real part of the spectrum is conditioned by $\cL$ and $\vth$. During the estimation process, it might therefore be that the condition $\delta\,\gamma \ll 1$ is violated.
% \item An interesting question (to be addressed as a next step) is what happens if the input projects on a fixed eigenvector of $\cL$. Do we observe something specific ?

% \item The rank of $\cM_F(\xi,\vth)$ is constrained by the rank of $\frac{\partial  \cL}{\partial \theta_p}$. We have assumed, from \eqref{eq:OrderEigenvalues}, that $\cL$ has full rank, but this does not tell us anything about its derivative. 
% %
% It is therefore possible, and even expected, that $\cM_F(\xi,\vth)$ is not invertible.
% \end{itemize}

% \subsection{Design}\label{Sec:Design}

% \subsubsection{Mathematical setting} \label{Sc:MathSetting}

% Here, we use the estimation \eqref{eq:theta_hat} of $\vth$, to propose a new set of stimuli. That is, we want to propose a set $\Set{\vu^{(1)}, \dots, \vu^{(k)}}$, distributed according to a probability density $\xi$, so as to minimize a criterion based on the empirical
% information matrix:
% %
% \begin{equation}\label{eq:InfoMatrix}
% \cM_F(\vu,\vth) \,=\, \sum_{k=1}^K \nu_k \, \pare{\frac{\partial \veta(\vth,\vu^{(k)})}{\partial \vth}}^T.\pare{\frac{\partial \veta(\vth,\vu^{(k])}}{\partial \vth}}.
% \end{equation}
% %
% This is an approximation of \eqref{eq:Fisher} where the exact density $\xi$ is replaced by the empirical frequencies $\nu_k$.
% The same computation as for \eqref{eq:Fisher} holds (eq. \eqref{eq:MF_exp}), where $\Exp{\xi}{}$  in eq. \eqref{eq:Delta} has to be replaced by the empirical average. Note that $K=\frac{P(P+1)}{2}$ samples are sufficient from \cite{pronzato:08} \BL{Je ne suis pas sur de cela. Pas sur que j'ai correctement compris.}

% Assume that we have obtained an estimation of $\vth$ using the criterion \eqref{eq:theta_hat}. From \cite{pronzato:08} we define the ellipsoid of confidence around $\hat{\vth^N_{LS}}$ as:
% %
% \begin{equation}\label{eq:Ellipsoid}
% \cR(\hat{\vth^N_{LS}},\alpha) \,=\, \Set{\vth \,|\, \pare{\vth \,-\,\hat{\vth^N_{LS}}}^T.\cM_F(\vu,\vth).\pare{\vth \,-\,\hat{\vth^N_{LS}}} \,\leq \, \frac{\chi^2_\alpha(P)}{N}},
% \end{equation}
% %
% where $\chi^2_\alpha(P)$ is a random variable, with a probability $\alpha$ to be exceeded by a random variable $z$ chi-square distributed with $P$ degrees of freedom: $\Prob{\chi^2_\alpha(P) \leq z} = \alpha$. We have:
% %
% $$
% \Prob{\vth \in \cR(\hat{\vth^N_{LS}},\alpha) } \,=\, \alpha.
% $$

% In this ellipsoid one attempts to minimize a criterion defined by a non linear function $\Phi$ of $\cM_F(\vu,\vth)$. Possible choices are:
% %
% \begin{itemize}
% \item \textbf{A optimal design.} Minimize $\Phi\pare{\cM_F}\,=\,Tr\pare{\cM_F^{-1}}$.
% \item \textbf{L optimal design.} Minimize $\Phi\pare{\cM_F}\,=\,Tr\pare{Q^T.\cM_F^{-1}.Q}$ for some weighting matrix $Q$.
% \item \textbf{E optimal design.} Minimize the longest axis of $\cR$, or, equivalently, maximising the minimum eigenvalue of $\cM_F$.
% \item \textbf{D optimal design.} Maximizing $\Phi\pare{\cM_F}\,=\,\det\cM_F$ , or, equivalently, minimising the volume of the ellipsoid $\cR(\hat{\vth^N_{LS}},\alpha)$.
% \end{itemize}

% In one dimension ($N=1$) there is a theorem stating that $D$ optimality is equivalent to minimizing the 
% distance $d_{\vth}$ defined below. Does this theorem still hold for $N>1$ ?. In dimension $N$ this distance would be:
% %
% \begin{equation}\label{eq:Distance_D}
% d_{\vth}\pare{\vu,\xi} \,=\,\pare{\frac{\partial \veta(\vth,\vu)}{\partial \vth}}^T.\cM_F^{-1}.\pare{\frac{\partial \veta(\vth,\vu)}{\partial \vth}}.
% \end{equation}

% \paragraph{Remarks.}

% \begin{itemize}
% \item $\cM_F(\vu,\vth)$ is a symmetric matrix, thus with real eigenvalues. Is there any connection between the eigenvalues of $\cM_F(\vu,\vth)$ and the eigenvalues of $\cL$ ?
% \item It can have zero eigenvalues due to possible degeneracies in the matrices $\frac{\partial  \cL}{\partial \theta_p}$ (see the remarks of the previous section).
% \end{itemize}

% \BL{
% \begin{itemize}
% \item How do we select $\alpha$ in \eqref{eq:Ellipsoid} ?
% \item Does the equivalence theorem between $D$ optimality and distance $d_{\vth}\pare{\vu,\xi}$ minimization hold in dimension $N$ ?
% \item It is not guaranteed that $\cM_F$ is invertible. Does it make sense to project the ellipsoid on $Ker^\perp \cM_F$ ?   
% \end{itemize}
% }

% \subsubsection{Algorithm}\label{Sec:Algorithm}

% We report here the Fedorov-Wynn algorithm, as presented in \cite{pronzato:08}.

% Let $\xi^{(m)}$ denote the distribution at iteration $m$ of the algorithm.

% \begin{enumerate}
% \item Choose $\xi^{(1)}$ non degenerate and $0<\epsilon \ll 1$.
% \item While $d_{\vth}\pare{\vu,\xi} > P \,+\, \epsilon$, 
% \begin{enumerate}
% \item Compute $\vu^{(m+1)}$ such that:
% %
% $$
% \vu^{(m+1)},=\,\argmax_{\vu \in \cU} d_{\vth}\pare{\vu,\xi}.
% $$
% \item Set $\alpha_m=\frac{(d_{\vth}\pare{\vu^{(m+1)},\xi}-p)}{\Set{p d_{\vth}\pare{\vu^{(m+1)},\xi}-1}}$.
% \item Set $\xi^{(m+1)}\,=\,(1-\alpha_k)\xi^{(m+1)} \,+\, \alpha_m \delta_{\vu^{(m+1)}}$.
% \item $m \to m+1$.
% \end{enumerate}
% \end{enumerate}

% \BL{A discuter.}






\section{LS, penalised LS and design of optimal experiments}

%\Lucinline{J'ai essay\'e d'utiliser les m\^emes notations que dans la note ``Parameter estimation in a linear model" de Bruno (ou du moins des notations aussi pr\`es que possible\ldots).}

We make observations $\yb(t_\ell,\ub)$ at times $t_\ell=(\ell-1)\delta$, $\ell=1,\ldots,H$. Each $\yb(t_\ell,\ub)$ is a $N$-dimensional vector (all vectors are column vectors) and corresponds to the model response to the piecewise constant input $\ub$ (an $H$-dimensional vector) corrupted by noise; that is,
\beq\label{model1}
\yb(t_\ell,\ub) = \etab(t_\ell,\bar\mtb,\ub)+\mveb(t_\ell) \,,
\eeq
with $\bar\mtb$ an unknown vector of $p$ parameters. We assume that the dynamic of the system is fast enough to make the errors $\mveb(t_\ell)$ and $\mveb(t_{\ell'})$ independent for $t_\ell\neq t_{\ell'}$ and denote $\Sigmab_\ell=\Ex\{\mveb(t_\ell)\mveb\TT(t_\ell)\}$, an $N \times N$ matrix, not necessarily diagonal.
%\Luccomment{Nous n'avons pas besoin de supposer que $\Sigmab_\ell$ est identique pour tous les $\ell$, ou est proportionnelle à la matrice identit\'e, ni m\^eme diagonale.}.

If the observations at different time instants were uncorrelated, with $\Sigmab_\ell$ known (not dependent of $\mtb$) for each $\ell$, 
the (weighted) LS criterion for the input $\ub$ would be: 
\bea
J(\mtb,\ub) = \sum_{\ell=1}^H \left[\yb(t_\ell,\ub) - \etab(t_\ell,\mtb,\ub)\right]\TT\Sigmab_\ell^{-1}\left[\yb(t_\ell,\ub) - \etab(t_\ell,\mtb,\ub)\right] \,.
\eea
When using all observations available for different inputs $\ub^{(k)}$, $k=1,\ldots,K$, we thus get the following criterion (after normalisation by $1/K$):
\bea
J_K(\mtb) = \frac1K\, \sum_{k=1}^K \sum_{\ell=1}^H \left[\yb(t_\ell,\ub^{(k)}) - \etab(t_\ell,\mtb,\ub^{(k)})\right]\TT\Sigmab_\ell^{-1}\left[\yb(t_\ell,\ub^{(k)}) - \etab(t_\ell,\mtb,\ub^{(k)})\right] \,.
\eea

In the following we account for the correlation between observations at different time instants (but assume that observations for different inputs $\ub^{(k)}$ correspond to different experiments and are uncorrelated). 
Denoting $\Yb(\ub)=\left[\yb\TT(t_1,\ub),\ldots,\yb\TT(t_H,\ub)\right]\TT$ and $\zetab(\mtb,\ub)=\left[\etab\TT(t_1,\mtb,\ub),\ldots,\etab\TT(t_H,\mtb,\ub)\right]\TT$, we get
\beq\label{LS-Sigma-known}
J_K(\mtb) = \frac1K\, \sum_{k=1}^K \left[\Yb(\ub^{(k)}) - \zetab(\mtb,\ub^{(k)})\right]\TT \Sigmab^{-1}\left[\Yb(\ub^{(k)}) - \zetab(\mtb,\ub^{(k)})\right] \,,
\eeq
where $\Sigmab$,
%\Luccomment{On pourrait prendre en compte la corr\'elation entre les $\mveb(t_\ell)$ et $\mveb(t_{\ell'})$ et utiliser une matrice pleine $\Sigmab$ --- mais $\Sigmab$ d\'ependrait alors de $\mtb$, nest-ce pas?} 
the $NH\times NH$ covariance matrix of the $\Yb(\ub^{(k)})$. %\B{C'est plutôt la matrice de covariance des $\epsilon$ non ?}\Luc{les $\yb(t_\ell,\ub)$ et les $\mveb(t_\ell)$ ne diffèrent que d'un terme déterministe, donc c'est indifférent}, does not depend on $\ub^{(k)}$.
%block diagonal matrix with $k$-th block given by $\Sigmab_\ell$.

%-----------------------------------------------
\subsection{Design of experiments: $\Sigmab$ is known (does not depend on $\bar\mtb$)}\label{S:Sigma-known}

Here we assume that $\Sigmab$ is known: it does not depend on $\bar\mtb$.

We consider the situation when $H$ is fixed but $K$ tends to infinity. Let $\xi$ denote a probability measure on the set $\SU\subset \mathds{R}^H$ of input vectors $\ub$. We assume that $\SU$ is compact and that
\begin{itemize}
  \item either $\xi$ has finite support and the empirical measure $\xi_{e,K}$ of the $\ub^{(k)}$ converges strongly to $\xi$ as $K\to\infty$ (i.e., $\lim_{K\to\infty}\xi_{e,K}(\{\ub\})=\xi(\{\ub\})$ for all $\ub\in\SU$),

  \item or the $\ub^{(k)}$ are independently sampled from $\xi$.
\end{itemize}

Then, under classical regularity assumptions on the functions $\etab(t,\cdot,\ub)$ and on the admissible set $\Theta$ for $\mtb$ ($\Theta$ is compact and $\bar\mtb\in \mathrm{int}(\Theta)$), the LS estimator $\hat\mtb^K$ that minimises $J_K(\mtb)$ with respect to $\mtb\in\Theta$ converges a.s.\ to $\bar\mtb$ when $K\to\infty$ and is asymptotically normal
\bea
\sqrt{K} (\hat\mtb^K-\bar\mtb) \rad \zb \sim \SN\left(\0b,\Mb^{-1}(\xi,\bar\mtb)\right)\,, \ K\to\infty\,,
\eea
where $\Mb(\xi,\mtb)$ denotes the information matrix for model \eqref{model1}:
\bea
\Mb(\xi,\mtb) = \int_\SU \frac{\mp\zetab\TT(\mt,\ub)}{\mp\mtb} \Sigmab^{-1} \frac{\mp\zetab(\mt,\ub)}{\mp\mtb\TT}\, \xi(\dd\ub) \,,
\eea
see, e.g., \cite[Chap.~3]{PP2013}, \cite[Chap.~1]{FedorovL2014}.

Optimal design (approximate theory) aims at determining an optimal probability measure $\xi^*$ on $\SU$ that minimises a suitable scalar function $\Phi(\cdot)$ of $\Mb(\xi,\bar\mtb)$. As $\bar\mtb$ is unknown, we consider locally optimal design which relies on $\Mb(\xi,\mtb^0)$ with $\mtb^0$ a nominal value for $\mtb$. Since each matrix
\bea
\Mb(\delta_\ub,\mtb)= [\mp\zetab\TT(\mt,\ub)/\mp\mtb] \Sigmab^{-1} [\mp\zetab(\mt,\ub)/\mp\mtb\TT]
\eea
has rank at least one, Caratheodory's theorem implies that for any $\xi$ on $\SU$ we can find a discrete $\xi'$ on $\SU$ supported on $m\leq p(p+1)/2+1$ points such that $\Mb(\xi',\mtb)=\Mb(\xi,\mtb)$. When $\Phi(\cdot)$ is monotone, $\Mb(\xi^*,\mtb)$ is on the boundary of the set $\{\Mb(\delta_\ub,\mtb):\ \ub\in\SU\}$ and $m\leq p(p+1)/2$ support points are enough. When $\rank[\Mb(\delta_\ub,\mtb)]=\dim(\mtb)$, the optimal design may correpond to a delta measure $\delta_\ub$ at a single input $\ub\in\SU$.

A typical design criterion corresponds to $D$-optimality, with $\Phi(\Mb)=-\log\det \Mb$. A multivariate extension, see for instance \cite{HarmanT2009}, of Kiefer-Wolfowitz (1960)  equivalence Theorem \cite{KieferW60} then states that an optimal $\xi^*$ for $\mtb^0$, a nominal value for the unknown true parameters  $\bar\mtb$, minimises $\max_{\ub\in\SU} d(\ub,\xi,\mtb^0)$, with, moreover, $\max_{\ub\in\SU} d(\ub,\xi^*,\mtb^0)=p$, where
\beq\label{D-function-Sigma-known}
d(\ub,\xi,\mtb) = \tr\left[ \frac{\mp\zetab\TT(\mt,\ub)}{\mp\mtb} \Sigmab^{-1} \frac{\mp\zetab(\mt,\ub)}{\mp\mtb\TT} \Mb^{-1}(\xi,\mtb) \right] \,.
\eeq

The vertex-direction algorithm of Fedorov (1972) \nocite{Fedorov72} %\Luccomment{Celui du \S 2.3.2 de la note de Bruno}
corresponds to a Frank-Wolfe algorithm for the maximization of $\log\det \Mb(\xi,\mtb^0)$. Starting from a measure $\xi^{(0)}$ such that $\Mb(\xi^{(0)},\mtb^0)$ is nonsingular, 
%\Luccomment{La non singularit\'e de $\Mb$ est essentielle.} 
at iteration $k\geq 0$ the algorithm chooses $\ub^{(k+1)}\in\Arg\max_{\ub\in\SU} d(\ub,\xi^{(k)},\mtb^0)$ and then sets
\bea
\xi^{(k+1)}=(1-\ma)\xi^{(k)}+\ma\delta_{\ub^{(k+1)}} \,.
\eea
The optimal value $\ma^*$ of the step size $\ma$ is known explicitly in the single response case, but must be determined numerically \B{How ?}\Luc{C'est un problème convexe de dimension 1, une méthode de type section dorée ou Fibonacci suffit} in the multi-response situation considered here.

\vsp
An alternative construction of an optimal design, which does not require differentiation of $\zetab(\mt,\ub)$ with respect to $\mtb$ (and thus no calculation of the information matrix) and does not rely on asymptotic considerations is proposed in \cite[Chap.~7]{PP2013} and \cite{PazmanP2014}. The criterion (to be maximised) is
\beq\label{phi-Sigma-known}
\phi(\xi) = \min_{\mtb\in\Theta} \frac{\int_\SU \left[\zetab(\mtb,\ub) - \zetab(\mtb^0,\ub)\right]\TT \Sigmab^{-1}\left[\zetab(\mtb,\ub) - \zetab(\mtb^0,\ub)\right]\,\xi(\dd\ub)}{\|\mtb-\mtb^0\|^2}\,,
\eeq
where $\Theta$ is a compact subset of $\mathds{R}^p$ containing admissible parameter vectors for the model. If we discretise $\Theta$ into a finite set $\Theta_Q=\{\mtb^{(1)},\ldots,\mtb^{(Q)}\}$ and consider discrete measures $\xi$ that set weights $w_i$ to $L$ support points $\ub^{(1)},\ldots,\ub^{(L)}$, the criterion becomes
\bea
\phi(\wb) = \min_{\mtb^{(j)}\in\Theta_Q} \frac{\sum_{i=1}^L w_i\, \left[\zetab(\mtb^{(j)},\ub^{(i)}) - \zetab(\mtb^0,\ub^{(i)})\right]\TT \Sigmab^{-1}\left[\zetab(\mtb^{(j)},\ub^{(i)}) - \zetab(\mtb^0,\ub^{(i)})\right]}{\|\mtb-\mtb^0\|^2}\,.
\eea
Its maximisation with respect to $\wb=(w_1,\ldots,w_L)\TT$ subject to $\sum_{i=1}^L w_i = 1$ corresponds to an LP problem:
\bea
&& \mbox{Maximise $t$ with respect to $t$ and $\wb\in\mathds{R}^L$} \B{N'est ce pas "Maximize $\phi$ ... ? }\\
%
&&\Luc{La formulation du pb de programmation lineaire dans $R^{L+1}$ me semble correcte}\\
%
&& \mbox{subject to } \sum_{i=1}^L w_i = 1 \mbox{ and the $Q$ additional linear constraints} \\
&& \hspace{1cm} \sum_{i=1}^L w_i \frac{\left[\zetab(\mtb^{(j)},\ub^{(i)}) - \zetab(\mtb^0,\ub^{(i)})\right]\TT \Sigmab^{-1}\left[\zetab(\mtb^{(j)},\ub^{(i)}) - \zetab(\mtb^0,\ub^{(i)})\right]}{\|\mtb-\mtb^0\|^2} \geq t \,, \\ && \hspace{9cm} j=1,\ldots,Q \,.
\eea
This is still a locally optimal design problem since we have kept $\mtb^0$ fixed, but we could also consider the minimum with respect to values $\mtb^{0,(k)}$ varying in a set $\Theta'_{Q'}$ (non intersecting $\Theta_Q$) --- at the expense of a larger number of linear constraints to be accounted for.

The construction could also be iterative. 
%\Luccomment{Je n'ai jamais essay\'e cette approche\ldots}
Suppose that at iteration $m$ we have used $n_m$ input vectors in $\SU$ with a design $\xi^{(m)}=\sum_{i=1}^L w_i^{(m)} \delta_{\ub^{(i)}}$ (i.e., each $\ub^{(i)}$ has been used $n_m\times w_i^{(m)}$ times), and that the estimated value of $\mtb$ is $\hat\mtb^{n_m}$. We want to choose $q$ additional inputs, so that $n_{m+1}$ will be equal to $n_m+q$. The LP problem to solve is then:
\bea
&& \mbox{Maximise $t$ with respect to $t$ and $\wb\in\mathds{R}^L$} \\
&& \mbox{subject to } \sum_{i=1}^L w_i = 1 \mbox{ and the $Q$ additional linear constraints} \\
&& \frac{1}{n_m+q} \sum_{i=1}^L (n_m w_i^{(m)} + q w_i) \\
&& \hspace{1cm} \times \frac{\left[\zetab(\mtb^{(j)},\ub^{(i)}) - \zetab(\hat\mtb^{n_m},\ub^{(i)})\right]\TT \Sigmab^{-1}\left[\zetab(\mtb^{(j)},\ub^{(i)}) - \zetab(\hat\mtb^{n_m},\ub^{(i)})\right]}{\|\mtb-\hat\mtb^{n_m}\|^2} \geq t \,, \\
&& \hspace{9cm} j=1,\ldots,Q \,.
\eea

%-----------------------------------------------
\subsection{Design of experiments: $\Sigmab=\Sigmab(\mtb)$ depends on $\mtb$}

\B{This is the case with our problem.}
A first possibility in that case is to use iteratively reweighted LS:
\begin{enumerate}
    \item set an initial value for $\Sigmab=\Sigmab^0$, for instance $\Sigmab(\mtb^0)$ with $\mtb^0$ a nominal value for $\mtb$;
    \item estimate $\mtb$ by minimising \eqref{LS-Sigma-known} with $\Sigmab=\Sigmab^0$; this gives $\hat\mtb^K_0$;
    \item repeat the estimation of $\mtb$ with now $\Sigmab=\Sigmab(\hat\mtb^K_0)$ in \eqref{LS-Sigma-known}; this gives $\hat\mtb^K_1$; substitute $\hat\mtb^K_1$ for $\hat\mtb^K_0$ and repeat, etc.
\end{enumerate}
From the asymptotic point of view (as $K\to\infty$), stopping at $\hat\mtb^K_1$ (i.e., two-stage LS) is enough, but this iteratively reweighted scheme can be continued until stage $m$ where $\hat\mtb^K_m \simeq \hat\mtb^K_{m-1}$.

Another possibility, which we develop below, is to use penalized LS: we assume that $\Yb(\ub)$ is normally distributed $\SN(\zetab(\bar\mtb,\ub),\Sigmab(\bar\mtb))$ \B{which is OK} and estimate $\mtb$ by maximum likelihood. The criterion to be minimised with respect to $\mtb$ is then
\beq
J_K(\mtb) &=& \frac1K\, \sum_{k=1}^K \left[\Yb(\ub^{(k)}) - \zetab(\mtb,\ub^{(k)})\right]\TT \Sigmab^{-1}(\mtb)\left[\Yb(\ub^{(k)}) - \zetab(\mtb,\ub^{(k)})\right] \nonumber \\
&& + \frac1K\,  \log\det \Sigmab(\mtb)\,, \label{penalised-LS}
\eeq
that is, due to the presence of the parameterised covariance $\Sigmab(\mtb)$, the penalty term $(1/K)\,  \log\det \Sigmab(\mtb)$ is added to the LS criterion \eqref{LS-Sigma-known}.
The information matrix is now 
\bea
\Mb(\xi,\mtb) = \int_\SU \Ex\left\{ \frac{\mp\log \varphi[\Yb(\ub)|\mtb]}{\mp\mtb} \frac{\mp\log \varphi[\Yb(\ub)|\mtb]}{\mp\mtb\TT}\right\} \, \xi(\dd\ub) \,,
\eea
where $\varphi[\Yb(\ub)|\mtb]$ denotes the p.d.f.\ of the normal distribution $\SN(\zetab(\mtb,\ub),\Sigmab(\mtb))$ and the expectation $\Ex\{\cdot\}$ is calculated for this distribution. Direct calculation gives, for all $i,j=1,\ldots,\dim(\mtb)$ % \B{NOt clear how you arrive to this, but OK},
\bea
\{\Mb(\xi,\mtb)\}_{i,j} &=& \int_\SU \frac{\mp\zetab\TT(\mt,\ub)}{\mp\mtb_i} \Sigmab^{-1}(\mtb) \frac{\mp\zetab(\mt,\ub)}{\mp\mtb_j}\, \xi(\dd\ub) \\
&& + \frac12\, \int_\SU \tr\left[\Sigmab^{-1}(\mtb)\frac{\mp\Sigmab(\mtb)}{\mp\mtb_i}\Sigmab^{-1}(\mtb)\frac{\mp\Sigmab(\mtb)}{\mp\mtb_j}\right]\,\xi(\dd\ub) \,.
\eea
Since $\Sigmab(\mtb)$ does not depend on $\ub$, we get 
\bea
\Mb(\xi,\mtb) &=& \int_\SU \frac{\mp\zetab\TT(\mt,\ub)}{\mp\mtb} \Sigmab^{-1}(\mtb) \frac{\mp\zetab(\mt,\ub)}{\mp\mtb\TT}\, \xi(\dd\ub) +  \Qb(\mtb) \,,
\eea
where 
\bea
\{\Qb\}_{i,j}=\frac12\, \tr\left[\Sigmab^{-1}(\mtb)\frac{\mp\Sigmab(\mtb)}{\mp\mtb_i}\Sigmab^{-1}(\mtb)\frac{\mp\Sigmab(\mtb)}{\mp\mtb_j}\right] \,.
\eea

The function $d(\cdot,\xi,\mtb)$ given by \eqref{D-function-Sigma-known}, which is used in a vertex-direction algorithm and to check the $D$-optimality of a given design, see Section~\ref{S:Sigma-known}, must be modified accordingly into
\bea
d(\ub,\xi,\mtb) = \tr\left\{\left[ \frac{\mp\zetab\TT(\mt,\ub)}{\mp\mtb} \Sigmab^{-1} \frac{\mp\zetab(\mt,\ub)}{\mp\mtb\TT}+ \Qb(\mtb)\right] \Mb^{-1}(\xi,\mtb) \right\} \,.
\eea
\B{Do we have the equivalent of  Kiefer-Wolfowitz theorem in this case too ?}\Luc{Rien ne change à part la matrice d'information.}

If $\bar\mtb$ denotes the true (unknown) value of $\mtb$, then, as $K\to\infty$, under the same assumptions about the regularity of the model and the distribution of the $\ub^{(k)}$ as in Section~\ref{S:Sigma-known}, $J_K(\mtb)$ tends to 
\bea
J_{\bar\mtb}(\mtb) &=& \int_\SU \left[\zetab(\mtb,\ub^{(k)}) - \zetab(\bar\mtb,\ub^{(k)})\right]\TT \Sigmab^{-1}(\mtb)\left[\zetab(\mtb,\ub^{(k)}) - \zetab(\bar\mtb,\ub^{(k)})\right] \, \xi(\dd\ub) \\
&& + \int_\SU \left\{\tr\left[\Sigmab(\bar\mtb)\Sigmab^{-1}(\mtb)\right] - \log \det\left[\Sigmab(\bar\mtb)\Sigmab^{-1}(\mtb)\right] \right\} \, \xi(\dd\ub) \\
&& + \int_\SU \log \det \Sigmab(\bar\mtb)\, \xi(\dd\ub) \\
&=& \int_\SU \left[\zetab(\mtb,\ub^{(k)}) - \zetab(\bar\mtb,\ub^{(k)})\right]\TT \Sigmab^{-1}(\mtb)\left[\zetab(\mtb,\ub^{(k)}) - \zetab(\bar\mtb,\ub^{(k)})\right] \, \xi(\dd\ub) \\
&& +\, \tr\left[\Sigmab(\bar\mtb)\Sigmab^{-1}(\mtb)\right] - \log \det\left[\Sigmab(\bar\mtb)\Sigmab^{-1}(\mtb)\right] + \log \det \Sigmab(\bar\mtb) \,.
\eea
and the convergence is uniform in $\mtb$. 
The criterion $J_{\bar\mtb}(\cdot)$ is minimum at $\mtb=\bar\mtb$, with $\mp J_{\bar\mtb}(\mtb)/\mp\mtb_{|\mtb=\bar\mtb}=\0b$.

Instead of maximising $\phi(\xi)$ given by \eqref{phi-Sigma-known}, we must now maximise
\beq\label{phi-penalised-LS}
\phi(\xi)= \min_{\mtb\in\Theta}  \frac{J_{\mtb^0}(\mtb)-J_{\mtb^0}(\mtb^0)}{\|\mtb-\mtb^0\|^2} = \min_{\mtb\in\Theta}  \frac{\int_\SU KL(\mtb||\mtb^0;\ub)\,\xi(\dd\ub)}{\|\mtb-\mtb^0\|^2} \,,
\eeq
where $KL(\mtb||\mtb^0;\ub)$
denotes the (nonsymmetric) Kullback-Leibler divergence between the normal distributions for models with parameters $\mtb$ and $\mtb^0$, with $\mtb^0$ a nominal value for the unknown true parameters  $\bar\mtb$:
\bea
KL(\mtb||\mtb^0;\ub) &=&\int \log\left\{\frac{\varphi[\Yb(\ub)|\mtb]}{\varphi[\Yb(\ub)|\mtb^0]}\right\}\, \varphi[\Yb(\ub)|\mtb]\,\dd\Yb \\
&=& \left[\zetab(\mtb,\ub) - \zetab(\mtb^0,\ub)\right]\TT \Sigmab^{-1}(\mtb)\left[\zetab(\mtb,\ub) - \zetab(\mtb^0,\ub)\right]  \\
&& +\, \tr\left[\Sigmab(\mtb^0)\Sigmab^{-1}(\mtb)\right] - \log \det\left[\Sigmab(\mtb^0)\Sigmab^{-1}(\mtb)\right] - \dim(\zetab)\,,
\eea
with $\dim(\zetab)=NH$. 

\B{Thus:
%
$$
\phi(\xi)\,=\, 
\min_{\mtb\in\Theta} \frac{1}{\|\mtb-\mtb^0\|^2} \,  \pare{ \begin{array}{ll}
&\int_\SU \left[\zetab(\mtb,\ub) - \zetab(\mtb^0,\ub)\right]\TT \Sigmab^{-1}(\mtb)\left[\zetab(\mtb,\ub) - \zetab(\mtb^0,\ub)\right] \, \xi(\dd\ub)\\
&+\, \tr\left[\Sigmab(\mtb^0)\Sigmab^{-1}(\mtb)\right] - \log \det\left[\Sigmab(\mtb^0)\Sigmab^{-1}(\mtb)\right] - \dim(\zetab).
\end{array}
%
}
$$
%
Denoting $\vs_i(\mtb)$ the normalized eigenvectors of $\Sigmab(\mtb)$, which are real and orthogonal, and $\sigma_i(\mtb)$ the corresponding real non negative eigenvalues
we can project $\zetab(\mtb,\ub) - \zetab(\mtb^0,\ub)$ on this basis, $\zetab(\mtb,\ub) - \zetab(\mtb^0,\ub) = \sum_{i=1}^{NH} \pi_i(\mtb,\ub) \vs_i(\mtb)$ and write:
%
$$
\int_\SU \left[\zetab(\mtb,\ub) - \zetab(\mtb^0,\ub)\right]\TT \Sigmab^{-1}(\mtb)\left[\zetab(\mtb,\ub) - \zetab(\mtb^0,\ub)\right] \, \xi(\dd\ub)
=\sum_{i=1}^{NH}  \frac{1}{\sigma_i(\mtb)} \int_\SU \pi_i^2(\mtb,\ub) \, \xi(\dd\ub).
$$
%
Note that this is the only term in the criterion $\phi$ which depends on $\xi$.
Also, we have:
%
$$
\log \det\left[\Sigmab(\mtb^0)\Sigmab^{-1}(\mtb)\right]  \,=\, \sum_{i=1}^{NH} \log \frac{\sigma_i(\mtb^0)}{\sigma_i(\mtb)}.
$$
%
Finally, since
$\Sigmab(\mtb^0) \,=\,\sum_{i=1}^{NH} \sigma_i(\mtb^0) \vs_i(\mtb^0) \vs_i(\mtb^0)^T$, $\Sigmab^{-1}(\mtb)\,=\,\sum_{j=1}^{NH} \sigma_j^{-1}(\mtb) \vs_j(\mtb) \vs_j(\mtb)^T$ we have:
%
$$
\Sigmab(\mtb^0)\Sigmab^{-1}(\mtb) \,=\, \sum_{i,j=1}^{NH} \frac{\sigma_i(\mtb^0)}{\sigma_j(\mtb)} \, \vs_i(\mtb^0) \, \underbrace{\vs_i(\mtb^0)^T \vs_j(\mtb)}_{<\vs_i(\mtb^0),\vs_j(\mtb)>} \vs_j(\mtb)^T
$$
%
so that:
%
$$
\pare{\Sigmab(\mtb^0)\Sigmab^{-1}(\mtb)}_{kl} \,=\, <\vs_k(\mtb),\pare{\Sigmab(\mtb^0)\Sigmab^{-1}(\mtb)}\vs_l(\mtb)>
$$
%
$$
=\, \sum_{i,j=1}^{NH} \frac{\sigma_i(\mtb^0)}{\sigma_j(\mtb)} \,<\vs_k(\mtb^0),\vs_i(\mtb)> \,<\vs_i(\mtb^0),\vs_j(\mtb)>  \,  \underbrace{<\vs_j(\mtb), \vs_l(\mtb)>}_{\delta_{jl}}
$$
%
$$
=\, \sum_{i=1}^{NH} \frac{\sigma_i(\mtb^0)}{\sigma_l(\mtb)} \, <\vs_k(\mtb^0),\vs_i(\mtb)> \,<\vs_i(\mtb^0),\vs_l(\mtb)> 
$$
and:
%
$$
\tr\left[\Sigmab(\mtb^0)\Sigmab^{-1}(\mtb)\right] \,=\, 
\sum_{i,k=1}^{NH} \frac{\sigma_i(\mtb^0)}{\sigma_k(\mtb)} \, <\vs_k(\mtb^0),\vs_i(\mtb)> \,<\vs_i(\mtb^0),\vs_k(\mtb)> 
$$
}

\bigskip

Direct calculation gives 
\bea
\frac{\mp^2 J_{\mtb^0}(\mtb)}{\mp\mtb\mp\mtb\TT}_{\big|\mtb=\mtb^0}= \frac{\int_\SU \mp^2 KL(\mtb||\mtb^0;\ub)\,\xi(\dd\ub)}{\mp\mtb\mp\mtb\TT}_{\big|\mtb=\mtb^0}=2\,\Mb(\xi,\mtb^0) \,.
\eea
Note that this forms a generalisation of the situation considered in Section~\ref{S:Sigma-known} where we assumed that $\Sigmab$ is known: when $\Sigmab$ does not depend on $\mtb$ (and thus $\Sigmab(\mtb^0)=\Sigmab(\mtb)$) we recover the same expressions for $\Mb(\xi,\mtb)$ and $\phi(\xi)$. 

%-----------------------------------------------
\subsection{Design of experiments: a (very!) small one-dimensional example}
% D:\Matlab\Design-m\B_Cessac\tst_dim2.m     

We take 
\bea
\cL=\left(\begin{array}{cc}
-1/\tau_A & - \omega^- \\
-\omega^+ & -1/\tau_B 
\end{array}\right)
\eea
in \eqref{eq:SD} with $\tau_A=0.1$, $\tau_B=0.05$, $\omega^-=0.1$, $\omega^+=0.2$; $\Fb(t)=\fb$, $\delta_{\epsilon}(t)$, with $\epsilon=0.1$ and $\fb=[f_1,f_2]'=[1,2]'$; $\ms_B^2=0.5$. We consider observations on $[0,0.2]$. Figure~\ref{F:y-and-eta} shows a particular realisation.

\begin{figure}[ht!]
\begin{center}
\includegraphics[width=.5\linewidth]{plots/y-and-eta}
\end{center}
\caption{\small $\yb(t)$ and $\etab(t)$.}\label{F:y-and-eta}
\end{figure}

To formulate a simple one-dimensional problem, we define $\mt=\tau_A+\tau_B$ and $u=f_2-f_1$ (with $\mt^0=\bar\mt = 0.15$ and $u=1$ in Figure~\ref{F:y-and-eta}). We observe $\yb(t)$ at 41 time instants equally spaced in $[0,0.2]$, for one single ``input signal": $K=1$ and $u=1$. The left panel of Figure~\ref{F:LS-criterion} presents the evolutions of three different estimation criteria as functions of $\mt\in\Theta=[0.1,0.2]$: $J_1(\mt)$ given by \eqref{LS-Sigma-known} with $\Sigmab=\Sigmab(\bar\mt)$ (an oracle LS criterion), $J_1(\mt)$ with $\Sigmab=\Sigmab(\mt)$ (a naive LS criterion that should not be used), and the penalised LS criterion \eqref{penalised-LS}. The estimated value of $\mt$ for each criterion is indicated. The right panel shows $J_1(\mt)/|\mt-\mt^0|^2$ as a function of $\mt$ when $\Sigmab=\Sigmab(\mt^0)$.

\begin{figure}[ht!]
\begin{center}
\includegraphics[width=.49\linewidth]{plots/LS-criterion}
\includegraphics[width=.49\linewidth]{plots/LS-criterion-divided-by-squared-variation-of-theta}
\end{center}
\caption{\small Left: $J_1(\mt)$, see \eqref{LS-Sigma-known} and \eqref{penalised-LS}. Right: $J_1(\mt)/|\mt-\mt^0|^2$ with $\Sigmab=\Sigmab(\mt^0)$.}\label{F:LS-criterion}
\end{figure}

The left panel of Figure~\ref{F:penalised-LS-criterion} shows $KL(\mt||\mt^0,u)/|\mt-\mt^0|^2$ as a function of $\mt$ for three different choices for $u$. The right panel shows $\phi(\xi)$ defined by \eqref{phi-penalised-LS} when $\xi$ is the delta measure at $u$, as a function of $u=f_2-f_1\in\SU=[-4,0]$. The optimal input (the opimal probability measure on $\SU$) is here $\delta_0$. This is consistent with the solution of the LP problem presented in Section~\ref{S:Sigma-known} obtained when we discretise the two sets $\Theta=[0.1,0.2]$ and $\SU=[-4,0]$. 

The situation is different when a delta measure yields a singular information matrix --- which is not the case here as $\mt$ is one-dimensional and furthermore $\Yb(u)$ has dimension $NH$. To obtain an interesting case we may think of a situation with either (\textit{i}) $\dim(\mtb)>NH$ or (\textit{ii}) when a given $\ub$ yields small values of $KL(\mtb||\mtb^0;\ub)$ when $\mtb$ lies in a subset $\Theta'(\ub)$ of $\Theta$ which depends on $\ub$. 

\begin{figure}[ht!]
\begin{center}
\includegraphics[width=.49\linewidth]{plots/KL-divided-by-squared-variation-of-theta}
\includegraphics[width=.49\linewidth]{plots/phi-single-u}
\end{center}
\caption{\small Left: $KL(\mt||\mt^0;u)/|\mt-\mt^0|^2$ for three different $u$. Right: $\phi(\delta_u)$ defined by \eqref{phi-penalised-LS} as a function of $u$.}\label{F:penalised-LS-criterion}
\end{figure}

%_______________________________________________
%_______________________________________________
%_______________________________________________
%_______________________________________________

\section{Appendix}\label{Sec:Appendix}

\subsection{Computation of the covariance}\label{Sec:ACov}

The covariance of  $\vX(t)$, the solution of \eqref{eq:SD} is defined by
%
$$
\cC(t,t') \deq \Exp{\vW}{\pare{\vX(t)-\mu(t)}.\pare{\vX(t')-\mu(t')}^T}.
$$
%
Using the general solution \eqref{eq:SolSD} and the definition of the expectation \eqref{eq:Exp_X} this is:
%
$$
\sigma_B^2 \, \int_{-\infty}^t e^{\cL(t-s)}.\underbrace{\Exp{\vW}{d\vW(s).d\vW^T(s')}}_{\delta(s-s')\,\cI}e^{\cLT(t'-s')} \,ds\,ds'
$$
%
$$
=\,\sigma_B^2 \, \int_{-\infty}^{t \wedge t'} e^{\cL(t-s)}.e^{\cLT(t'-s)} \,ds.
$$
%
Assuming $t\leq t'$ this gives:
%
$$
\cC(t,t')  \,=\, \sigma_B^2 \, \int_{-\infty}^t e^{\cL(t-s)}.e^{\cLT(t'-s)} \, ds, 
$$
%
corresponding to \eqref{eq:Cov_X} in the main text.
Using the diagonal form of $\cL$ gives:
%
$$
\cC(t,t')  \,=\, \sigma_B^2 \, \int_{-\infty}^t \cP.e^{\Lambda(t-s)}.\cP^{-1}. \pare{\cP.e^{\Lambda(t'-s)}.\cP^{-1}}^T\, ds
$$
%
$$
=\, \sigma_B^2 \, \int_{-\infty}^t  \cP.e^{\Lambda(t-s)}.\underbrace{\cP^{-1}.\pare{\cP^{-1}}^T}_{\cG}.e^{\Lambda(t'-s)}.\cP^T.
$$

The matrix $\cG\,=\,\cP^{-1}.\pare{\cP^{-1}}^T$ can also be written in terms of the left eigenvectors of $\cL$ (the rows of $\cP^{-1}$):
%
$$
\cG \,=\, \pare{\vw^T_1, \dots, \vw^T_N}.\vect{\vw_1\\\dots\\\vw_N}.
$$
%
This is a positive, definite, matrix with entries
%
$g_{mn}\,=\,\vw^T_m.\vw_n\,=\,<\vw_m,\vw_n>$.
%
This defines therefore a metric. 

We note that $\cD(t,t') \equiv \cP^{-1}.\cC(t,t').\pare{\cP^{-1}}^T$
has entries:
%
$$
\cD_{mn}(t,t') \,=\, \sigma_B^2 \,g_{mn} \,  \int_{-\infty}^t e^{\lambda_m(t-s)}\,e^{\lambda_n(t'-s)}\, ds
\,=\, \sigma_B^2 \,g_{mn} \, e^{\lambda_m \, t \,+\, \lambda_n \, t'} \int_{-\infty}^t 
e^{-\,\pare{\lambda_m \,+\, \lambda_n}s} \, ds
$$
$$
=\, - \, \sigma_B^2 \, \frac{g_{mn}}{\lambda_m \,+\, \lambda_n} \, e^{\lambda_m \, t \,+\, \lambda_n \, t'} \, 
\bra{e^{-\,\pare{\lambda_m \,+\, \lambda_n}s}}_{-\infty}^t
$$ 
%
$$
\cD_{mn}(t,t') \,=\,  
- \, \sigma_B^2 \, \frac{g_{mn}}{\lambda_m \,+\, \lambda_n} \, e^{\lambda_n \, \pare{t'-t}},
$$
using that all eigenvalues are stable. Defining the matrix $\cA$ with entries $A_{mn}=\frac{g_{mn}}{\lambda_m \,+\, \lambda_n}$ we have
$\cD(t,t') \,=\,- \, \sigma_B^2 \cA.e^{\Lambda(t-t')}$.

It is also useful to compute the inverse of $\cD(t,t')$, $\cD^{-1}(t,t')\,=\ \cP^T.\cC^{-1}(t,t').\cP$, which is:
%
\begin{equation}\label{eq:Dmoins1}
\cD^{-1}(t,t') \,=\, - \, \sigma_B^{-2} \, e^{-\Lambda(t-t')}.\cA^{-1}.
\end{equation} 

\subsection{Computation of $\frac{\partial \cI^{l,r}}{\partial \theta_p}$}\label{Sec:dIdtheta}

From \eqref{eq:Ilr}, 
$\cI^{l,r}\,=\,\int_{\delta r}^{\delta (r+1)} e^{\cL\pare{l \delta - s}} \,ds$
%
so, $\frac{\partial \cI^{l,r}}{\partial \theta_p}\,=\, \int_{\delta r}^{\delta (r+1)} \frac{\partial e^{\cL\pare{l \delta - s}}}{\partial \theta_p} \,ds$. In general, $\frac{\partial  \cL}{\partial \theta_p}$ does not commute with $\cL$  so one cannot write $\frac{\partial e^{\cL\pare{l \delta - s}}}{\partial \theta_p}\,=\,\pare{l \delta - s} \frac{\partial \cL}{\partial \theta_p}.e^{\cL\pare{l \delta - s}}$. Instead, using the series definition of $e^{\cL\pare{t - s}}$ one can show that:
%
$$
\frac{\partial e^{\cL\pare{t - s}}}{\partial \theta_p} \,=\, \sum_{n=1}^{+\infty} \frac{\pare{t - s}^n}{n!} \pare{\sum_{m=0}^{n-1} \cL^m.\frac{\partial  \cL}{\partial \theta_p}.\cL^{n-(m+1)}}.
$$
%
Thus:
%
$$
\cP^{-1}.\frac{\partial e^{\cL\pare{t - s}}}{\partial \theta_p}.\cP\,=\, \sum_{n=1}^{+\infty} \frac{\pare{t - s}^n}{n!} \pare{\sum_{m=0}^{n-1} \Lambda^m.\cD_p.\Lambda^{n-(m+1)}}
$$
%
with:
%
\begin{equation}\label{eq:Dp}
\cD_p\,=\,\cP^{-1}.\frac{\partial  \cL}{\partial \theta_p}.\cP.
\end{equation}

We have therefore:
%
$$
\cP^{-1}.\frac{\partial \cI^{l,r}}{\partial \theta_p}.\cP \,=\, \sum_{n=1}^{+\infty} \frac{1}{n!} \pare{\sum_{m=0}^{n-1} \Lambda^m.\cD_p.\Lambda^{n-(m+1)}} \int_{\delta r}^{\delta (r+1)} \pare{l \delta - s}^n\,ds 
$$
%
$$
=\,\sum_{n=1}^{+\infty} \frac{\delta^{n+1}}{(n+1)!} \pare{\sum_{m=0}^{n-1} \Lambda^m.\cD_p.\Lambda^{n-(m+1)}} \,A_{l,r,n+1},
$$
%
with $A_{l,r,n+1} \,=\,\bra{(l-r)^{n+1} \,-\,(l-r-1)^{n+1}} \,=\, \sum_{a=1}^{n+1} \binom{n+1}{a} (-1)^{1+a} \, (l-r)^{n+1-a}.$
%
Setting $\cG^{l,r}_p=\sum_{n=1}^{+\infty} \frac{\delta^{n+1}}{(n+1)!} \pare{\sum_{m=0}^{n-1} \Lambda^m.\cD_p.\Lambda^{n-(m+1)}} \,A_{l,r,n+1}$
This gives:
%
\begin{equation}\label{eq:dIdtheta}
\frac{\partial \cI^{l,r}}{\partial \theta_p} \,=\,
\cP.\cG^{l,r}_p.\cP^{-1}.
\end{equation}

Note that we can rewrite $\cG^{l,r}_p$ in the form:
%
$$
\cG^{l,r}_p\,=\,\delta^{2} \,\sum_{n=1}^{+\infty} \frac{1}{(n+1)!} \pare{\sum_{m=0}^{n-1} \pare{\delta\Lambda}^m.\cD_p.\pare{\delta\Lambda}^{n-(m+1)}} \,A_{l,r,n+1},
$$
%
letting appear the dimensionless matrix $\delta\Lambda$ controlling the convergence rate of the series.

\subsection{Computation of the Fisher matrix in terms of the eigen-elements of $\cL$}\label{Sec:ComputeFisher}

Using \eqref{eq:DiscX}-\eqref{eq:eta} we have:
%
$$
\eta_q(\vth,\vu)\,=\, \sum_{j=1}^N \, \sum_{r=0}^{l-1} \, \cI^{l,r}_{ij}\,F_j(r\delta), \, q\,=\, i+lN.
$$
%
As $F_j(r\delta)$ does not depend on $\vth$ we have:
%
$$
\frac{\partial \eta_q(\vth,\vu)}{\partial \theta_p}
\equiv E_{qp}
\,=\,\sum_{j=1}^N \, \sum_{r=0}^{l-1} \, \frac{\partial \cI^{l,r}_{ij}}{\partial \theta_p}\,F_j(r\delta),
$$
%
corresponding to a $Q \times P$ matrix. Then:
%
$$
\bra{\pare{\frac{\partial \veta(\vth,\vu)}{\partial \vth}}^T.\pare{\frac{\partial \veta(\vth,\vu)}{\partial \vth}}}_{pp'}
\,=\, \sum_{q=1}^Q E_{pq}\,E_{qp'}
$$
$$=\, \sum_{i=1}^N \sum_{l=1}^H E_{p,i+lN}\,E_{i+lN,p'}
\,=\, \sum_{i=1}^N \sum_{l=1}^H \pare{\sum_{j=1}^N \, \sum_{r=0}^{l-1} \, \frac{\partial \cI^{l,r}_{ij}}{\partial \theta_p}\,F_j(r\delta)}
\pare{\sum_{j'=1}^N \, \sum_{r'=0}^{l-1} \, \frac{\partial \cI^{l,r'}_{ij'}}{\partial \theta_{p'}}\,F_{j'}(r'\delta)},
$$
%
$$
=\,\ \sum_{l=1}^H \sum_{r,r'=0}^{l-1} \, \sum_{i=1}^N \,\sum_{j,j'=1}^N \,\frac{\partial \cI^{l,r}_{ij}}{\partial \theta_p}\,F_j(r\delta)\,F_{j'}(r'\delta)\, \frac{\partial \cI^{l,r'}_{ij'}}{\partial \theta_{p'}}. 
$$

Therefore $\cM_F(\xi,\vth)$ as entries:
%
%
This is eq. \eqref{eq:MF_exp} in the main text, with
%
$$
\Delta_{jj'}(r,r') \,=\, \Exp{\xi}{F_j(r\delta)\,F_{j'}(r'\delta) }.
$$


Note that \eqref{eq:MF_exp} can also be written as:
%
\begin{equation}\label{eq:Fisher_Trace}
\cM_{pp'}(\xi,\vth)\,=\,  \sum_{l=1}^H \sum_{r,r'=0}^{l-1} \, Tr\bra{\frac{\partial \cI^{l,r}}{\partial \theta_p}.\Delta(r,r').\frac{\partial \cI^{l,r'}}{\partial \theta_{p'}}},
\end{equation}
%
where $l,r,r'$ are time indices and where $\Delta(r,r')\,=\,\Exp{\xi}{\vF(r\delta).\vF^T(r'\delta)}$. 

It is useful to write this matrix in terms of $\vf(r\delta)\,=\,\cP^{-1}.\vF(r\delta)$, the projection of the input in the eigenspace of $\cL$.
Using \eqref{eq:dIdtheta} we have:
%
$$
\frac{\partial \cI^{l,r}}{\partial \theta_p}.\Delta(r,r').\frac{\partial \cI^{l,r'}}{\partial \theta_{p'}}
\,=\, \cP.\cG^{l,r}_p.\cP^{-1}.\Exp{\xi}{\vF(r\delta).\vF^T(r'\delta)}.\cP.\cG^{l,r'}_p.\cP^{-1}
$$
%
$$
=\, \cP.\cG^{l,r}_p.\underbrace{\Exp{\xi}{\vf(r\delta).\vf^T(r'\delta)}}_  {\cC_f(r,r')}.\cP^T.\cP.\cG^{l,r'}_p.\cP^{-1},
$$

Thus:
%
\begin{equation}\label{eq:Fisher_Propre}
\cM_{pp'}(\xi,\vth)\,=\,  \sum_{l=1}^H \sum_{r,r'=0}^{l-1} \, Tr\bra{
\cP.\cG^{l,r}_p.\cC_f(r,r').\cP^T.\cP.\cG^{l,r'}_p.\cP^{-1}
}.
\end{equation}

\B{Que se passe-t'il si l'input est projeté sur un seul vecteur propre ?}

In this decomposition, the inputs are  projected onto a fixed basis, only the coefficients of the projection are random, governed by $\xi$.
They impact the Fisher matrix via the correlation matrix $\cC_f(r,r')$ (it's a correlation if the coefficients are centered). 



%{\ludo 


\subsection{Continuous-discrete model (Ludovic)}

Let us denote:
%
\begin{equation}
    \cW(t)=\int_{-\infty}^t e^{\cL(t-s)}d W(s).
\end{equation}
%
Naturally, $\Exp{}\cW\equiv0$, and we already have $\cC(t,t')$ denoting the covariance of $(\cW(t),\cW(t'))$. Following Gardiner, $\cC(t,t)=\cC(0,0)$ in the stationary limit, is solution of the (well-posed) Sylvester equation\footnote{It is even a Lyapunov equation.}: $\cL\cC(t,t)+\cC(t,t)\cL^{T}=-\cI_N$. Furthermore, if $t>t'$,
$\cC(t,t')=e^{\cL(t-t')}\cC(t,t)$ ($t'<t$ recovered via $\cC(t,t')=\cC(t',t)^T$).

We use these expressions to discretize \eqref{eq:SolSDGen} under the assumption that $F$ is constant over $[t_l,t_{l+1})$\footnote{One big advantage is that we do not even need to assume that $F$ is constant, just that the input doesn't change. For instance it would work the same assuming $F$ is a sinusoid.}.
\begin{equation}
    \vX(t_{l+1})=e^{\delta\cL}\vX(t_{l})+\cL^{-1}(e^{\delta\cL}-\cI_N)F(t_l)+\cW(t_{l+1})-\cW(t_{l})
\end{equation}
Letting $\cM=e^{\delta\cL}$, we obtain 
\begin{equation}
\begin{aligned}
    \vX(t_{l})=&\sum_{n=0}^{l-1}\cM^{l-n-1}\cL^{-1}(\cM-\cI_N)F(t_n)
    \\
    &+\sum_{n=-\infty}^{l}\cM^{l-n}\cW(t_{n})-\sum_{n=-\infty}^{l}\cM^{l-n}\cW(t_{n-1})
    \\
    =&(\cM-\cI_N)\sum_{n=0}^{l-1}\cM^{l-n-1}\cL^{-1}F(t_n)
    \\
    &+W(t_l)+(\cM-\cI_N)\sum_{n=-\infty}^{l-1}\cM^{l-n-1}\cW(t_{n}).
\end{aligned}
\end{equation}


\begin{lemma}
The Sylvester (matrix) equation $X-\cM X \cM^T=Y$ always admits a unique solution for any $Y$.
\end{lemma}
\begin{proof}
I admit the well known fact that the Syvester equation $AX-XB=C$ admits a unique solution if and only if the spectrum of $A$ and $B$ are disjoint %\B{Don't we have conditions on the rank of $A$ and $B$ ? e.g. $A$ and $B$ have full rank ?}\L{C'est impliqué par les conditions de spectre: si les deux ont un noyau non trivial, les spectres ne sont pas disjoints.}. 
Since $\cM=e^{\delta\cL}$, it is invertible and $X-\cM X \cM^T=Y$ is equivalent to $\cM^{-1}X-X \cM^T=\cM^{-1}Y$, so the characterization applies %\B{What do you mean by "the characterization" ? That $\cM^{-1}$ and $cM^T$ have disjoint spectrum ?}
%\L{Je veux dire: le résultat sur l'existence de solutions passe de la forme $AX-XB=C$ à $X-AXB=C$ dans ces conditions }. 
Recall it was assumed on the spectrum of $\cL$:
\begin{equation}\label{E:spectrumL}
\re{\lambda_N} \leq \re{\lambda_{N-1}} \leq \dots \leq \re{\lambda_1} < \gamma<0.    
\end{equation}
The spectrum of $\cM$ is given by $\sigma(\cM)=\left(e^{\lambda_i}\right)_{1\leq i\leq n}$. In particular, $\sigma(\cM^T)=\sigma(\cM)$ must belong to the interior of the disk of radius $1$ in $\mathbb{C}$ by \eqref{E:spectrumL}. But then $\sigma(\cM^{-1})=\left(e^{-\lambda_i}\right)_{1\leq i\leq n}$ belongs to the interior of the complement of disk of radius $1$ in $\mathbb{C}$. Therefore  $\sigma(\cM^{-1})\cap \sigma(\cM^T)=\emptyset$ and the statement holds.
\end{proof}



Let $\varepsilon(t_l)=W(t_l)+(\cM-\cI_N)\sum_{n=-\infty}^{l-1}\cM^{l-n-1}\cW(t_{n})$. Then $\Exp{}\varepsilon\equiv0$ and for $l>l'$, we have 
\begin{multline}
    \mathrm{Cov}(\varepsilon(t_l),\varepsilon(t_{l'}))
    =
    \cC(t_l,t_{l'})+
    \left\{
    \sum_{n'=-\infty}^{l'-1}\cC(t_l,t_{n'}){(\cM^T)}^{l'-n'-1}
    \right\}(\cM^T-\cI_N)
    \\
    +
    (\cM-\cI_N)\sum_{n=-\infty}^{l-1}{\cM}^{l-n-1}\cC(t_n,t_{l'})
    \\
    +
    (\cM-\cI_N)
    \left\{
        \sum_{n'=-\infty}^{l'-1}\sum_{n=-\infty}^{l-1}{\cM}^{l-n-1}\cC(t_n,t_{n'}){(\cM^T)}^{l'-n'-1}
    \right\}
    (\cM^T-\cI_N).
\end{multline}

With $\cC(0,0):=\cC_0$, we have
$$
\cC(t_{n},t_{n'})=
\begin{cases}
\cM^{n-n'}\cC_0& \text{ if } n\geq n',
\\
\cC_0 {(\cM^{T})}^{n'-n}& \text{ if } n< n'.
\end{cases}
$$
Consider
$$
\cC_{1}:=\sum_{n'=-\infty}^{l'-1}\cM^{l'-n'-1}\cC_0{(\cM^T)}^{l'-n-1}=\sum_{k=0}^{\infty}\cM^{k}\cC_0{(\cM^T)}^{k}.
$$
The matrix $\cC_{1}$ satisfies the Sylvester equation $\cC_1-\cM\cC_1\cM^T=\cC_0$. We have
\begin{equation}\label{E:term2}
\sum_{n'=-\infty}^{l'-1}\cC(t_l,t_{n'}){(\cM^T)}^{l'-n'-1}
=
\sum_{n'=-\infty}^{l'-1}\cM^{l-n'}\cC_0{(\cM^T)}^{l'-n'-1}
=
\cM^{l-l'+1}\cC_{1}
,
\end{equation}
and
\begin{equation}\label{E:term3}
\begin{aligned}
\sum_{n=-\infty}^{l-1}{\cM}^{l-n-1}\cC(t_n,t_{l'})
=
&\sum_{n=-\infty}^{l'-1}{\cM}^{l-n-1}\cC_0 {(\cM^{T})}^{l'-n}
+
\sum_{n=l'}^{l-1}{\cM}^{(l-n-1)+(n-l')}\cC_0
\\
=&
\cM^{l-l'}\cC_{1}\cM^T+(l-l')\cM^{l-l'-1}\cC_{0}.
\end{aligned}
\end{equation}
Finally, since $n'\leq l'-1<l-1$, we can divide the sum
\begin{equation}
\begin{aligned}
\sum_{n=-\infty}^{l-1}{\cM}^{l-n-1}\cC(t_n,t_{n'}){(\cM^T)}^{l'-n'-1}
    =&\sum_{n=-\infty}^{n'-1}{\cM}^{l-n-1}\cC(t_n,t_{n'}){(\cM^T)}^{l'-n'-1}
    \\
    &+
    \sum_{n=n'}^{l-1}{\cM}^{l-n-1}\cC(t_n,t_{n'}){(\cM^T)}^{l'-n'-1}
\end{aligned}
\end{equation}
so that 
\begin{equation}
\begin{aligned}
    \sum_{n=-\infty}^{n'-1}{\cM}^{l-n-1}\cC(t_n,t_{n'}){(\cM^T)}^{l'-n'-1}
    =&
    \sum_{n=-\infty}^{n'-1}{\cM}^{l-n-1}\cC_0{(\cM^T)}^{n'-n}{(\cM^T)}^{l'-n'-1}
    \\
    =&
    \sum_{n=-\infty}^{n'-1}{\cM}^{l-n-1}\cC_0{(\cM^T)}^{l'-n-1}
    \\
    =&
    {\cM}^{l-n'}\cC_1{(\cM^T)}^{l'-n'}
\end{aligned}
\end{equation}
and
\begin{equation}
\begin{aligned}
    \sum_{n=n'}^{l-1}{\cM}^{l-n-1}\cC(t_n,t_{n'}){(\cM^T)}^{l'-n'-1}
    =&
    \sum_{n=n'}^{l-1}{\cM}^{l-n-1}{\cM}^{n-n'}\cC_0{(\cM^T)}^{l'-n'-1}
    \\
    =&
    \sum_{n=n'}^{l-1}{\cM}^{l-n'-1}\cC_0{(\cM^T)}^{l'-n'-1}
    \\
    =&
    (l-n'){\cM}^{l-n'-1}\cC_0{(\cM^T)}^{l'-n'-1}.
\end{aligned}
\end{equation}
Then
\begin{equation}
\begin{aligned}
\sum_{n'=-\infty}^{l'-1}
    \left\{
        \sum_{n=-\infty}^{l-1}{\cM}^{l-n-1}\cC(t_n,t_{n'}){(\cM^T)}^{l'-n'-1}
    \right\} 
=&
\sum_{n'=-\infty}^{l'-1}{\cM}^{l-n'}\cC_1{(\cM^T)}^{l'-n'}
\\
&+
\sum_{n'=-\infty}^{l'-1}(l-n'){\cM}^{l-n'-1}\cC_0{(\cM^T)}^{l'-n'-1}
\end{aligned}
\end{equation}
Denote
$$
\cC_{2}:=\sum_{n'=-\infty}^{l'-1}\cM^{l'-n'-1}\cC_1{(\cM^T)}^{l'-n'-1}=\sum_{k=0}^{\infty}\cM^{k}\cC_1{(\cM^T)}^{k},
$$
solution of the Sylvester equation $\cC_2-\cM\cC_2\cM^T=\cC_1$,
and 
$$
\cC_{1}':=\sum_{n'=-\infty}^{l'-1}(l'-n'-1){\cM}^{l'-n'-1}\cC_0{(\cM^T)}^{l'-n'-1}=\sum_{k=1}^{\infty}k \cM^{k}\cC_0{(\cM^T)}^{k}.
$$
Some sum manipulation shows that $\cC_1'$ is solution to the Sylvester equation 
$\cC_1'-\cM\cC_1'\cM^T=\cC_1-\cC_0=\cM\cC_1\cM^T$.

Eventually,
\begin{equation}\label{E:term4}
\sum_{n'=-\infty}^{l'-1}
    \left\{
        \sum_{n=-\infty}^{l-1}{\cM}^{l-n-1}\cC(t_n,t_{n'}){(\cM^T)}^{l'-n'-1}
    \right\} 
    =\cM^{l-l'+1} \cC_2 \cM^T+(l-l'+1)\cM^{l-l'}\cC_1+\cM^{l-l'}\cC_1'.
\end{equation}

The conclusion of the computations is
\begin{multline}
    \mathrm{Cov}(\varepsilon(t_l),\varepsilon(t_{l'}))
    =
    \cM^{l-l'}\cC_0+\cM^{l-l'+1}\cC_{1}(\cM^T-\cI_N)
    \\
    +
    (\cM-\cI_N)\cM^{l-l'}\cC_{1}\cM^T+(l-l')(\cM-\cI_N)\cM^{l-l'-1}\cC_{0}
    \\
    +
    (\cM-\cI_N)\cM^{l-l'}
    \left\{
        \cM \cC_2 \cM^T+(l-l'+1)\cC_1+\cC_1'
    \right\}
    (\cM^T-\cI_N).
\end{multline}
This can be slightly simplified. Recall that $\cC_1-\cM\cC_1\cM^T=\cC_0$, $\cC_2-\cM\cC_2\cM^T=\cC_1$ and $\cC_1'-\cM\cC_1'\cM^T=\cC_1-\cC_0$. This yields that $\cC'_1=\cC_2-\cC_1$. Then 
$$
\cM \cC_2 \cM^T+(l-l'+1)\cC_1+\cC_1'=\cM \cC_2 \cM^T+\cC_2+(l-l')\cC_1.
$$

\B{It seems to me that these expression are quite more tractable by projecting on the eigenbasis of $\cM$.}
\L{Il me semble qu'en grande dimension au contraire, cela coute moins cher tout en étant plus précis. Toutefois la réduction idéale dans ce contexte c'est la décomposition de Schur, à voir si elle aide à écrire les calculs.}


\paragraph{Let me redo the computations when $l=l'$.}

\begin{multline}
    \mathrm{Cov}(\varepsilon(t_l),\varepsilon(t_{l}))
    =
    \cC(t_l,t_{l})+
    \left\{
    \sum_{n'=-\infty}^{l-1}\cC(t_l,t_{n'}){(\cM^T)}^{l-n'-1}
    \right\}(\cM^T-\cI_N)
    \\
    +
    (\cM-\cI_N)\sum_{n=-\infty}^{l-1}{\cM}^{l-n-1}\cC(t_n,t_{l})
    \\
    +
    (\cM-\cI_N)
    \left\{
        \sum_{n'=-\infty}^{l-1}\sum_{n=-\infty}^{l-1}{\cM}^{l-n-1}\cC(t_n,t_{n'}){(\cM^T)}^{l'-n'-1}
    \right\}
    (\cM^T-\cI_N).
\end{multline}
From \eqref{E:term2},
\begin{equation}\label{E:term2bis}
\sum_{n'=-\infty}^{l-1}\cC(t_l,t_{n'}){(\cM^T)}^{l-n'-1}
=
\cM\cC_{1}
,
\end{equation}
and from \eqref{E:term3},
\begin{equation}\label{E:term3bis}
\sum_{n=-\infty}^{l-1}{\cM}^{l-n-1}\cC(t_n,t_{l})
=
\cC_{1}\cM^T.
\end{equation}
Likewise, from \eqref{E:term4},
\begin{equation}\label{E:term4bis}
\sum_{n'=-\infty}^{l-1}
    \left\{
        \sum_{n=-\infty}^{l-1}{\cM}^{l-n-1}\cC(t_n,t_{n'}){(\cM^T)}^{l'-n'-1}
    \right\} 
    =\cM \cC_2 \cM^T+\cC_1+\cC_1'.
\end{equation}
In conclusion,
\begin{multline}\label{E:concl}
    \mathrm{Cov}(\varepsilon(t_l),\varepsilon(t_{l}))
    =
    \cC_0+\cM\cC_{1}(\cM^T-\cI_N)
    +
    (\cM-\cI_N)\cC_{1}\cM^T
    \\
    +
    (\cM-\cI_N)
    \left\{
        \cM \cC_2 \cM^T+\cC_2
    \right\}
    (\cM^T-\cI_N).
\end{multline}
Interestingly, for $l\geq l$,
\begin{equation}\label{E:covllp}
\mathrm{Cov}(\varepsilon(t_l),\varepsilon(t_{l'}))
=
\cM^{l-l'}\left(
\mathrm{Cov}(\varepsilon(t_l),\varepsilon(t_{l}))+
(l-l')\cC_1
\right)
\end{equation}
(This can probably be proved without computing both and then noticing the formula \B{which formula ?}\L{Désolé pour la rédaction brouillonne. Si le but est de prouver la formule de $\mathrm{Cov}(\varepsilon(t_l),\varepsilon(t_{l'}))$, il est plus efficace de calculer $\mathrm{Cov}(\varepsilon(t_l),\varepsilon(t_{l}))$ puis de prouver que \eqref{E:covllp} vaut, plutôt que de calculer les deux indépendamment.}, to simplify the proofs.)
Now I want to remove as many terms as possible from the expression~\eqref{E:concl}, using the fact that $\cC_1,\cC_2$ satisfy Sylvester equations. For instance, $\cM\cC_2\cM^T=\cC_2-\cC_1$:
\begin{equation}
    \begin{aligned}
    \mathrm{Cov}(\varepsilon(t_l),\varepsilon(t_{l}))
    =&
    \cC_0+\cM\cC_{1}(\cM^T-\cI_N)
    +
    (\cM-\cI_N)\cC_{1}\cM^T
    \\
    &+
    (\cM-\cI_N)
    \left\{
        2 \cC_2 -\cC_1
    \right\}
    (\cM^T-\cI_N).
    \\
    =&
    \cC_0+\cM\cC_{1}\cM^T-\cM\cC_{1}+\cM\cC_{1}\cM^T-\cC_{1}\cM^T
    \\
    &+
    2\cM
        \cC_2
    \cM^T
    -
    2\cM
        \cC_2
    -
    2\cC_2
    \cM^T
    +
    2\cC_2
    \\
    &
    -
    \cM
       \cC_1
    \cM^T
    +\cM
       \cC_1
    +
       \cC_1
    \cM^T
    -
       \cC_1
    \\
    =&
    \cC_0+\cM\cC_{1}\cM^T-\cC_1
    \\
    &+
    2\cM
        \cC_2
    \cM^T
    -
    2\cM
        \cC_2
    -
    2\cC_2
    \cM^T
    +
    2\cC_2.
    \end{aligned}
\end{equation}
Then $\cM\cC_{1}\cM^T-\cC_1=-\cC_0$ yields 
\begin{equation}
    \begin{aligned}
    \mathrm{Cov}(\varepsilon(t_l),\varepsilon(t_{l}))
    =&
    2\cM
        \cC_2
    \cM^T
    -
    2\cM
        \cC_2
    -
    2\cC_2
    \cM^T
    +
    2\cC_2
    \\
    =&
    4\cC_2-2\cC_1-2\cM\cC_2-2\cC_2
    \cM^T.
    \end{aligned}
\end{equation}

\L{L'idée ici serait de laisser à une libriaire le soin de calculer les matrices $\cC_i$ de façon efficace, puis de se ramener à des produits matriciels et exponentiations rapides}

\paragraph{An interesting formula for the derivative of the matrix exponential.}

\begin{lemma}
For any $t>0$, 
\begin{equation}
\partial_{\theta_p} e^{t\cL}
=
\begin{pmatrix}
    \cI_N & 0_N
\end{pmatrix}
\exp\left( t
\begin{pmatrix}
    \cL &\partial_{\theta_p} \cL
    \\
    0&\cL
\end{pmatrix}
\right)
\begin{pmatrix}
    0_N 
    \\
    \cI_N 
\end{pmatrix}
\end{equation}
\end{lemma}
\begin{proof}
Let $\cA=\begin{pmatrix}
    \cL &\partial_{\theta_p} \cL
    \\
    0&\cL
    \end{pmatrix}$.
We prove the identity:
\begin{equation}
    \exp t
    \cA
    =
    \begin{pmatrix}
    e^{t\cL} &\partial_{\theta_p} e^{t\cL}
    \\
    0&e^{t\cL}
    \end{pmatrix}.
\end{equation}
First, we can check that: 
\begin{equation}
    \exp t
    \cA
    =
    \begin{pmatrix}
    e^{t\cL} &F(t)
    \\
    0&e^{t\cL}
    \end{pmatrix}
\end{equation}
with $F(t)=\int_0^t e^{(t-s)\cL}\partial_{\theta_p} \cL e^{s\cL}ds$.
This follows from the fact that $F(0)=0$ and $F'(t)=\cL F(t)+\partial_{\theta_p} \cL e^{t\cL}$.
Then we can recognize that $F(t)=\partial_{\theta_p}e^{t\cL}$, for instance from \cite[Equation (2.1)]{Wilcox1967}.
\end{proof}

\begin{corollary}
Let $\delta>0$, let $\cM=e^{\delta \cL}$. Let 
$$
\bar{\cM}_p=\exp\left( \delta
\begin{pmatrix}
    \cL &\partial_{\theta_p} \cL
    \\
    0&\cL
\end{pmatrix}
\right)
$$
For every positive integer $l$,
$$
\partial_{\theta_p}(\cM^l)
=
\begin{pmatrix}
    \cI_N & 0_N
\end{pmatrix}
{\left(\bar{\cM}_p\right)}^l
\begin{pmatrix}
    0_N 
    \\
    \cI_N 
\end{pmatrix}
$$
\end{corollary}

\B{Numerically, we may use a power expansion of the form:
$$
\bar{\cM}_p
= \cI_{2N} \,+\, \delta\, \begin{pmatrix}
    \cL &\partial_{\theta_p} \cL
    \\
    0&\cL
\end{pmatrix} \,+\, \frac{\delta^2}{2}
\begin{pmatrix}
    \cL^2 &\cL\,\partial_{\theta_p} \cL \,+\, \partial_{\theta_p} \cL
    \cL \\
    0&\cL^2
\end{pmatrix} \,+\, \dots
$$
%
}
%Let $\cM_p=\partial_{\theta_p}\cM$. If we set $\eta(\theta,t_l)=\vX(t)-\varepsilon(t_l)$

\L{La philosphie est d'éviter cette approximation: on utilise une librairie pour calculer rapidement $\cM$, $\bar{\cM}_p$, etc... ensuite tout se calcule via des produits matriciels, ce qui est peu couteux.}

\B{Comment calculer une exponentielle de matrice rapidement sans passer par la diagonalisation ?}
%}



\bibliographystyle{abbrv}
\bibliography{odyssee,Luc}

\end{document}